\documentclass[11pt]{article}

% -----------------------------------------------------------------------------
% Preamble supplied by the authors.  The paper body is intentionally contained
% in this file; in particular, there are no \input or \subfile commands.
% -----------------------------------------------------------------------------
\usepackage{fullpage}
\usepackage[T1]{fontenc}
\usepackage[utf8]{inputenc}
\usepackage{lmodern} % Scalable fonts for the compiled conference manuscript.
\usepackage{xcolor}
\definecolor{red}{HTML}{F88379}
\usepackage[colorlinks=true]{hyperref}
\hypersetup{
  linkcolor=blue,
  citecolor=orange,
  filecolor=purple,
  urlcolor=orange,
  menucolor=purple,
  runcolor=purple,
  pdftitle={Consensus Clustering below Two: Online Recourse and Fully Dynamic Algorithms},
  pdfsubject={Online consensus clustering with recourse and fully dynamic consensus clustering under vertex insertions and deletions}
}
\usepackage{amsmath,amssymb,amsthm,amsfonts,graphicx,color,subcaption,
  enumerate,bm,array,mathtools}
\usepackage[capitalise,nameinlink]{cleveref}
\usepackage{multicol,multirow}
\usepackage{thmtools,thm-restate}
\usepackage{ragged2e}
\usepackage{lineno}
\sloppy
\usepackage{framed}
\usepackage[framemethod=tikz]{mdframed}
% Comment this out for no line numbers.
% \linenumbers

% The supplied preamble uses algorithm2e.  All algorithms below are typeset as
% framed figures so that this source also compiles in minimal TeX installations.
\IfFileExists{algorithm2e.sty}{
  \usepackage[linesnumbered,ruled,vlined]{algorithm2e}
}{}

\graphicspath{{graphics/}}
\crefname{claim}{Claim}{Claims}
\crefname{property}{Property}{Properties}
\crefname{algocf}{Algorithm}{Algorithms}
\Crefname{algocf}{Algorithm}{Algorithms}

% Overleaf has the bibliography setup used in the supplied preamble.  The
% fallback keeps the same source compilable in a minimal local TeX installation.
\newif\ifusebiblatex
\IfFileExists{biblatex.sty}{
  \usebiblatextrue
  \IfFileExists{trad-alpha.bbx}{
    \usepackage[style=trad-alpha,natbib=true,maxcitenames=4]{biblatex}
  }{
    \usepackage[style=alphabetic,natbib=true,maxcitenames=4]{biblatex}
  }
  \addbibresource{ref.bib}
}{
  \usebiblatexfalse
}
\usepackage{dirtytalk}

% A dedicated algorithm float keeps the pseudocode presentation independent
% of the particular algorithm2e version installed by the compiler.
\DeclareCaptionType{paperalgorithm}[Algorithm][List of Algorithms]
\crefname{paperalgorithm}{Algorithm}{Algorithms}
\Crefname{paperalgorithm}{Algorithm}{Algorithms}

\makeatletter
\g@addto@macro\bfseries{\boldmath}
\makeatother

\newtheorem{theorem}{Theorem}[section]
\newtheorem{lemma}[theorem]{Lemma}
\newtheorem{claim}[theorem]{Claim}
\newtheorem{fact}[theorem]{Fact}
\newtheorem{corollary}[theorem]{Corollary}
\newtheorem{observation}[theorem]{Observation}
\newtheorem{proposition}[theorem]{Proposition}
\newtheorem{property}[theorem]{Property}
\theoremstyle{definition}
\newtheorem{definition}[theorem]{Definition}
\newtheorem{axiom}[theorem]{Axiom}
\theoremstyle{remark}
\newtheorem{remark}[theorem]{Remark}
\newtheorem*{remark*}{Remark}

% -----------------------------------------------------------------------------
% Notation.  These commands replace the external macros file used by the old
% project, keeping the manuscript self-contained apart from ref.bib.
% -----------------------------------------------------------------------------
\DeclareMathOperator*{\argmin}{arg\,min}
\newcommand{\A}{\mathcal{A}}
\newcommand{\C}{\mathcal{C}}
\newcommand{\F}{\mathcal{F}}
\newcommand{\Hh}{\mathcal{H}}
\newcommand{\I}{\mathcal{I}}
\newcommand{\Pp}{\mathcal{P}}
\newcommand{\Qq}{\mathcal{Q}}
\newcommand{\Ss}{\mathcal{S}}
\newcommand{\Y}{\mathcal{Y}}
\newcommand{\OPT}{\operatorname{OPT}}
\newcommand{\cost}{\operatorname{cost}}
\newcommand{\err}{\operatorname{err}}
\newcommand{\dist}{\operatorname{dist}}
\newcommand{\Agree}{\ensuremath{\mathsf{Agree}}}
\newcommand{\Pivot}{\ensuremath{\mathsf{Pivot}}}
\newcommand{\eps}{\varepsilon}
\newcommand{\1}{\mathbf{1}}
\newcommand{\E}{\mathbb{E}}
\newcommand{\Prb}{\mathbb{P}}
\newcommand{\Vtwo}{\binom{V}{2}}
\newcommand{\polylog}{\operatorname{polylog}}
\newcommand{\wtO}{\widetilde{O}}
\newcommand{\restr}[2]{#1|_{#2}}
\newcommand{\same}{\mathsf{same}}
\newcommand{\sep}{\mathsf{sep}}
\allowdisplaybreaks

\title{Consensus Clustering below Two:\\
Online Recourse and Fully Dynamic Algorithms}
\author{Anonymous Authors}
\date{}

\begin{document}
\pagenumbering{roman}
\maketitle
\thispagestyle{empty}

\begin{abstract}
Consensus clustering asks for a partition minimizing its average pairwise
disagreement with $m$ input partitions.  We study this metric-median problem
when the ground set changes over time, with two objectives: limiting changes
to previously assigned vertices and limiting the time needed to maintain a
solution.  We obtain approximation guarantees strictly below $2$ in both
settings.

We answer both questions through a common sampled-triple principle.  For three
input partitions, let $H$ be their majority signed graph.  If $F$ is any
reference partition and $\Delta(F)$ is its total triangle-inequality slack on
the triple, then the number of majority signs on which $H$ disagrees with $F$
is at most $\Delta(F)/2$.  This half-slack bound, derived from an exact overlap
identity, lets a correlation-clustering algorithm on $H$ produce a new median
candidate that is substantially better than merely returning a sampled input.

In the online setting, vertices arrive one at a time.  We maintain sampled
input candidates and persistent triple-majority candidates produced by the
deterministic Online Agreement algorithm.  If $\eta$ denotes the latter's
constant approximation factor, then for every
$0<\eps<\frac{2}{3\eta+5}$ we obtain a randomized algorithm with competitive ratio
\[
  2-\frac{2}{3\eta+5}+\eps
\]
with probability $1-\zeta$ simultaneously over every prefix of an oblivious
arrival sequence.  The polylogarithmic per-vertex recourse bound holds for
every realization of the initial samples.  Monotone history scores and a
multiplicative rebuilding rule convert the pointwise sampled-triple phenomenon
into a guarantee stable over the entire arrival history.

For fully dynamic consensus clustering, we maintain a
$(13/7+\eps)$-approximation under arbitrary vertex insertions and deletions.
For an oblivious sequence of at most $T$ updates and at most $N$ active
vertices, the guarantee holds simultaneously after every update with
probability $1-\zeta$.  Writing
$\delta=\min\{1/14,\eps/40\}$ and
$L=\log(8(N+1)(T+1)/(\delta\zeta))$, the worst-case time per update is
$O(\delta^{-3}(n_t+1)L^2)$ and the maintained space is
$O(\delta^{-2}(n_t+1)L)$ words, where $n_t$ is the current number of
vertices.  These bounds assume access to sampled coordinates of an arriving
vertex's input-label vector; scanning the full vector adds $O(m)$ insertion
time.  A variant withstands an adaptive adversary by using fresh samples and
random access to the full current input collection.

The dynamic algorithm rebuilds from a near-linear-time offline sampler.  Its
local routine runs \Pivot{} implicitly on a triple-majority graph by grouping
vertices according to pairs of cluster labels, avoiding its potentially
quadratic edge set.  A separate evaluation sample selects a good candidate.
The dynamic result makes no low-recourse claim; the online result makes no
near-linear update-time claim.
\end{abstract}

\newpage
\bigskip
{\small\tableofcontents}
\bigskip
\thispagestyle{empty}

\newpage
\pagenumbering{arabic}

% =============================================================================
\section{Introduction}
\label{sec:introduction}
% =============================================================================

Consensus clustering is the problem of aggregating several, possibly
inconsistent, partitions of the same data set into a single representative
partition.  The input partitions may be produced by different clustering
algorithms, by different parameter choices, by multiple views of the data, or
by independent executions of a randomized procedure.  Rather than declaring
one input authoritative, consensus clustering seeks the partition that best
captures their common structure.

The standard objective is pairwise.  Two partitions disagree on a pair of
elements if one puts the pair in a common cluster and the other separates it.
Given input partitions $P_1,\ldots,P_m$ of a ground set $V$, the cost of a
partition $C$ is
\[
  \cost_{\Pp}(C)=\frac{1}{m}\sum_{i=1}^{m}\dist(C,P_i),
\]
where $\dist$ counts pair disagreements.  This formulation exposes two useful
viewpoints.  First, consensus clustering is a $1$-median problem in the metric
space of partitions.  Second, if $p_{uv}$ is the fraction of inputs that
separate $u$ and $v$, then a solution pays $p_{uv}$ when it joins the pair and
$1-p_{uv}$ when it separates the pair.  It is therefore a highly structured
weighted correlation-clustering instance.

The problem is computationally difficult even for a constant number of input
partitions.  Early work established its general hardness~\cite{FilkovSkiena2003},
and it is APX-hard already for three inputs~\cite{BDDJ2008}.  At the same time, its metric
form immediately gives a simple benchmark.  If an input partition is selected
uniformly at random, the triangle inequality implies an expected
$2(1-1/m)$-approximation.  Thus, the value $2$ is both easy to approach and a
natural barrier to cross.

The seminal work of Ailon, Charikar, and Newman~\cite{ACN2008} crossed this
barrier by exploiting the correlation-clustering representation.  Combining a
majority-based pivot solution with an input candidate gives an expected
$11/7$-approximation; combining an LP-based candidate with an input candidate
gives a $4/3$-approximation.  These classical guarantees remain the principal
benchmarks for the general consensus objective.  Their original algorithms
are designed around pairwise information: explicitly constructing the
weighted instance or its majority graph takes $\Theta(n^2)$ pair evaluations.
Cordner and Kollios~\cite{CordnerKollios2023} showed how to
implement majority-based methods on the input labels in $O(nm)$ time and
studied sampling a prescribed number of input partitions.  Their work makes
clear that dependence on the number of inputs is a second bottleneck: when
$m$ is large, even reading the explicit $mn$-entry representation is costly.

This bottleneck is not removed by fast algorithms for ordinary unweighted
correlation clustering.  Sparse--dense decompositions and parallel or
near-linear techniques have produced increasingly efficient algorithms for
that problem~\cite{AW2022,CLMNPT2021,CALP2024}, but the signed graph supplied
to them must still be made available.  In a consensus instance,
determining the majority sign of every pair from all $m$ inputs is itself the
expensive operation.  There is also an objective mismatch: after subtracting
the unavoidable pairwise baseline, consensus clustering weights a majority
mistake by the pair's margin $|1-2p_{uv}|$, whereas an ordinary unweighted
correlation-clustering routine charges every majority mistake equally.  Thus,
neither materializing the graph nor simply discarding its heterogeneous
margins provides the desired input-sampling algorithm.  Our offline question
is therefore the following:

\begin{quote}
Can one obtain a worst-case, high-probability approximation strictly below
$2$ in time near-linear in $n$, independent of $m$, while inspecting only
polylogarithmically many input partitions?
\end{quote}

\paragraph{The online frontier.}
The same factor-$2$ barrier takes a different form when the ground set evolves.
Suppose vertices arrive one at a time.  Upon the arrival of a vertex, the
algorithm learns its cluster label in every input partition, and it must output
a consensus partition of all vertices seen so far.  Recomputing an offline
solution after every arrival gives no stability guarantee: one new vertex may
change the assignments of many old vertices.  At the other extreme, following
one fixed input partition has zero recourse but again approaches factor $2$.

We measure stability by \emph{per-vertex recourse}: the number of times an old
vertex changes its persistent cluster identifier.  For ordinary unweighted
correlation clustering, Cohen-Addad, Lattanzi, Maggiori, and
Parotsidis~\cite{cohen2022online} gave a deterministic constant-competitive
algorithm with $O(\log n)$ worst-case recourse per vertex.  Its guarantee
concerns the unweighted disagreement objective on a signed graph; it does not
directly provide a sub-$2$ approximation to the weighted consensus objective.
The online question is thus:

\begin{quote}
Can a fixed gap below $2$ hold simultaneously at every prefix, with high
probability, while every vertex changes clusters only polylogarithmically many
times?
\end{quote}

\paragraph{The dynamic frontier.}
When vertices may also be deleted, the active data set need not grow and its
optimum cost need not be monotone.  A deletion can remove precisely the pairs
responsible for a previously large objective value.  A history-based
certificate suitable for arrivals can therefore become arbitrarily larger
than the current optimum.  For this model, we ask for a different guarantee:

\begin{quote}
Can a sub-$2$ approximation be maintained after arbitrary vertex insertions
and deletions in $n_t\polylog(NT)$ time per update, where $n_t$ is the current
number of active vertices, while inspecting only polylogarithmically many
input partitions?
\end{quote}

An explicit answer already contains $n_t$ cluster labels, so a near-linear
bound is a natural target when the full current clustering is to be produced.
Here recourse and update time are separate resources: writing all labels is
affordable within the target time, but their changes need not be infrequent.
Our solution to this question uses a fast snapshot algorithm and a sampling
argument that remains valid when the active vertex set shrinks.

\subsection{Our results}
\label{subsec:intro-results}

We obtain affirmative answers in the models stated below.  The online and
dynamic algorithms optimize different resources, but their analyses share a
combinatorial core: if sampled inputs themselves do not contain a good median,
then a random triple has small triangle-inequality slack with nonnegligible
probability, and the majority relation of that triple is close to the unknown
optimum.

For the online result, we use the following black-box parameters from Online
Agreement~\cite{cohen2022online}.  Let $\eta\ge 1$ be an absolute constant such
that the algorithm is $\eta$-competitive at every prefix of an unweighted
complete signed graph, and let $q_{\mathrm A}$ be an absolute constant such
that each vertex changes its identifier at most
$q_{\mathrm A}\lceil\log_2 n\rceil$ times.  The published analysis leaves
$\eta$ unspecified; throughout, $\eta$ may be any certified upper bound on
that constant, after the internal agreement parameter is fixed to a sufficiently
small absolute value.

For a prefix $V_t$, let $\cost_t(C)$ be the average distance from $C$ to the
$m$ restricted inputs and let $\OPT_t=\min_C\cost_t(C)$.  An instance is
oblivious when the full partitions and arrival order are fixed independently
of the algorithm's initial sampling.

\begin{theorem}[Online consensus clustering]
\label{thm:online-main}
Let $\zeta\in(0,1)$ and
$\eps\in\bigl(0,\frac{2}{3\eta+5}\bigr)$.  Against every oblivious instance of
online consensus clustering whose horizon $n$ is known in advance, there is a
randomized online algorithm such that,
with probability at least $1-\zeta$, simultaneously for every prefix $t$,
\[
  \cost_t(C_t)
  \le
  \left(2-\frac{2}{3\eta+5}+\eps\right)\OPT_t.
\]
For every outcome of the initial sampling, every vertex has recourse
\[
  O_{\eta,q_{\mathrm A}}\!\left(
    \eps^{-2}
    \log\frac{2n}{\zeta}
    \log\frac{2mn}{\zeta}
  \right).
\]
\end{theorem}

The recourse bound holds deterministically for every realization of the
initial sampling.  The competitive ratio is independent of the number of
input partitions, while the recourse depends only logarithmically on $m$.
For every fixed $\eps$ and fixed black-box constants, the theorem gives a
constant gap below $2$ and polylogarithmic per-vertex recourse.  The result
optimizes competitive ratio and recourse; it does not claim a near-linear
update-time implementation.

\paragraph{Fully dynamic maintenance.}
There are $m$ indexed input partitions on a changing active set $V_t$.
Inserting a vertex supplies its cluster label in each input; deleting a vertex
restricts every input to the surviving vertices.  Existing labels and pair
relations among survivors are unchanged.  Let $n_t=|V_t|$, and assume known
bounds $N\ge1$ on $n_t$ and $T\ge1$ on the number of updates.  The normalized
objective $\cost_t$ and optimum $\OPT_t$ are defined on the current set.
The sampled-coordinate model permits querying selected entries of an arriving
vertex's label vector in constant time per entry.  An oblivious update
sequence, including the labels supplied on insertions and the vertices chosen
for deletion, is fixed independently of the algorithm's randomness.

\begin{theorem}[Fully dynamic consensus clustering: oblivious updates]
\label{thm:dynamic-main}
Let $\eps>0$ and $\zeta\in(0,1)$, and define
\begin{equation}
  \delta:=\min\{1/14,\eps/40\},\qquad
  L:=\log\frac{8(N+1)(T+1)}{\delta\zeta}.
  \label{eq:dynamic-parameters}
\end{equation}
There is a randomized algorithm in the sampled-coordinate model such that,
for every oblivious sequence of at most $T$ vertex insertions and deletions,
with probability at least $1-\zeta$, simultaneously after every update,
\[
  \cost_t(C_t)\le\left(\frac{13}{7}+\eps\right)\OPT_t.
\]
After an update leaving $n_t$ active vertices, it takes worst-case time
$O(\delta^{-3}(n_t+1)L^2)$, including materializing the output, and uses
$O(\delta^{-2}(n_t+1)L)$ words of maintained space.  It retains labels only
for a fixed sample of $O(\delta^{-2}L)$ input indices.  These resource bounds
hold for every execution.  If reading the entire incoming label vector is
charged, an insertion incurs an additional $O(m)$ time.
\end{theorem}

For constant $\eps$, a polynomial update horizon $T=N^{O(1)}$, and inverse
polynomial failure probability, the bounds become
$O_{\eps}((n_t+1)\log^2(N+1))$ update time and
$O_{\eps}((n_t+1)\log(N+1))$ space.  Deletions may be chosen adversarially
in advance; no distributional assumption on which vertex is deleted is needed.
The $+1$ accounts for empty and singleton snapshots and the stored samples.
The theorem imposes no recourse bound on surviving vertices.

\paragraph{Adaptive updates with full input access.}
An adaptive adversary may select the next update after seeing the previous
outputs, but before seeing the fresh private randomness for that update.
To handle this stronger adversary, we change the access model: after each
update, the algorithm can read any selected input partition on the entire
current active set.  This requires a database containing all current inputs;
storing only a fixed set of sampled coordinates does not provide this access.

\begin{theorem}[Fully dynamic consensus clustering: adaptive updates]
\label{thm:adaptive-main}
With the parameters in \eqref{eq:dynamic-parameters}, suppose the full current
collection of $m$ input-label arrays is available for random access after
every update.  There is a randomized algorithm against the adaptive adversary
described above that, with probability at least $1-\zeta$, simultaneously
after at most $T$ vertex insertions and deletions satisfies
\[
  \cost_t(C_t)\le\left(\frac{13}{7}+\eps\right)\OPT_t.
\]
Its worst-case clustering computation per update is
$O(\delta^{-3}(n_t+1)L^2)$, using
$O(\delta^{-2}(n_t+1)L)$ words of working space in addition to the input
database.  These bounds exclude maintaining the full input database.  If it
is stored locally with one length-$m$ label vector per active vertex, it uses
$O(m(n_t+1))$ words and can be
maintained with an additional $O(m+\log(N+1))$ time per insertion or deletion.
\end{theorem}

Fresh samples are necessary for our adaptive proof: conditioning on the past
may correlate a previously fixed sample with the next instance.  By contrast,
the oblivious algorithm samples its input indices once and stores only those
coordinates.  Thus the adversary distinction has a concrete space and
input-access consequence, even though the clustering computation has the same
asymptotic bound in the two theorems.

\paragraph{The snapshot algorithm.}
Both dynamic theorems follow by rebuilding with an offline algorithm whose
running time is already near-linear.  We state this supporting result
separately.  It assumes random access to an indexed collection of input
partitions, each represented by an $n$-entry label array.  Sampling one input
index costs $O(1)$, and reading the corresponding partition costs $O(n)$.

\begin{theorem}[Near-linear-time offline consensus clustering]
\label{thm:offline-main}
For every constant $\eps>0$, there is a randomized algorithm that, given
label-array representations of $m$ partitions of an $n$-element ground set,
returns with high probability a partition $C$ satisfying, for
$\OPT:=\min_F\cost_{\Pp}(F)$,
\[
  \cost_{\Pp}(C)\le\left(\frac{13}{7}+\eps\right)\OPT.
\]
The algorithm runs in $\wtO_{\eps}(n)$ time, uses $\wtO_{\eps}(n)$ working
space, and reads only $O_{\eps}(\log n)$ input partitions.  In particular, its running
time is independent of $m$ apart from constant-time access to uniformly
sampled indices.
\end{theorem}

The approximation succeeds with high probability; the time, working-space,
and input-access bounds hold for every execution.  The theorem complements
the $O(nm)$ majority-based implementation of
Cordner--Kollios~\cite{CordnerKollios2023} and their expected approximation
tradeoff for sampling a prescribed number of inputs.  Our constant is weaker
than $11/7$, but it comes with a high-probability guarantee and a worst-case
bound of only a polylogarithmic number of inspected inputs; the running time is
near-linear in the number of vertices even when $m$ is very large.

\subsection{Technical overview}
\label{subsec:techniques}

We next explain the common approximation argument, the history control needed
online, and the probability and access accounting needed for dynamic updates.

\paragraph{Why sampled inputs alone stop at two.}
Fix an optimal median $F^*$.  The random-input argument gives
\[
  \E_i[\cost_{\Pp}(P_i)]
  =\frac{1}{m^2}\sum_{i,j}\dist(P_i,P_j)
  \le 2\left(1-\frac1m\right)\OPT.
\]
This expectation does not imply that a constant fraction of the inputs has
cost at most $(2-\gamma)\OPT$ for any fixed $\gamma>0$.  Indeed, the fraction
can vanish as $m$ grows.  Therefore, sampling more input medoids, by itself,
cannot provide a fixed gap below $2$.

For an explicit example, take $m\ge3$ disjoint vertex pairs and let $F^*$
cluster each pair separately.  Let $P_i$ be obtained from $F^*$ by splitting
only pair $i$.  The optimum normalized cost is $1$: each of the $m$ pairs
contributes $1/m$, and joining any vertices from different pairs only increases
the cost.  Every input has cost $2-2/m$.  Thus even inspecting all input
partitions and returning their best member cannot yield a fixed gap below
$2$ for this family.  A successful algorithm must be able to synthesize a new
partition.

\paragraph{Slack and majority fitting.}
For each input $P_i$, let
\[
  I_i=E^+(P_i)\mathbin{\triangle}E^+(F^*)
\]
be the set of pairs on which $P_i$ and $F^*$ disagree.  For two inputs,
\[
  \dist(P_i,P_j)=|I_i|+|I_j|-2|I_i\cap I_j|.
\]
Consequently, their triangle-inequality slack through $F^*$ is exactly
\[
  \Delta_{ij}(F^*)
  :=\dist(F^*,P_i)+\dist(F^*,P_j)-\dist(P_i,P_j)
  =2|I_i\cap I_j|.
\]
Now sample three inputs and label a pair positive exactly when at least two of
the three co-cluster it.  The majority relation disagrees with $F^*$ only if
at least two inputs err on the pair.  If $B$ denotes this majority-error set,
then
\[
  B=(I_1\cap I_2)\cup(I_1\cap I_3)\cup(I_2\cap I_3),
  \qquad
  |B|\le\frac12\sum_{i<j}\Delta_{ij}(F^*).
\]
The factor $1/2$ is essential: it improves the local constant in the offline
analysis from $4$ to $2$ and correspondingly enlarges the gap obtained in the
online analysis.

\paragraph{The local-to-global dichotomy.}
A universal averaging lemma for metric $1$-median says the following.  Fix
$\alpha>0$.  Either at least a small fraction $\rho$ of the inputs are already
$(2-\alpha)$-approximate medians, or a random constant-size sample has expected
slack at most $O(\alpha+\rho)$ times $\OPT$.  In the first case, sampling input
partitions finds a good candidate.  In the second, repeatedly sampling triples
finds a majority instance close to $F^*$.

For a sampled triple $Q$, write
$\Delta_Q(F^*):=\sum_{\{i,j\}\in\binom Q2}\Delta_{ij}(F^*)$.

For the offline algorithm, \Pivot{} has expected error at most three times the
optimum correlation-clustering error of the triple-majority graph.  Passing
through the majority relation gives
\[
  \E[\dist(F^*,\Pivot(H_Q))]\le 2\Delta_Q(F^*).
\]
Thus the local constant is $C=2$.  The generic balance
\[
  \alpha=\frac{1}{C\binom{3}{2}+1}
\]
gives $\alpha=1/7$ and the ratio $2-\alpha=13/7$.

\paragraph{Making the offline local solver linear.}
The triple-majority graph may be dense, so it is never constructed.  A vertex
$v$ is represented by the three cluster labels $(c_1(v),c_2(v),c_3(v))$.
Two vertices are adjacent exactly when at least two coordinates agree.  The
closed neighborhood of $v$ is therefore the union of the three buckets keyed
by
\[
  (c_1(v),c_2(v)),\qquad
  (c_1(v),c_3(v)),\qquad
  (c_2(v),c_3(v)).
\]
We store the active vertices of each bucket in a doubly linked list.  A pivot
step enumerates only the union of its three current buckets and immediately
deletes every enumerated vertex.  Since every vertex is deleted once, the
entire \Pivot{} execution takes $O(n)$ time and space.  A separate contingency-
table identity computes the distance between two label-array partitions in
$O(n)$ time, allowing all candidates to be evaluated on a fresh polylogarithmic
sample of inputs.

\paragraph{History stability in the online setting.}
A pointwise good candidate is insufficient online.  If the algorithm repeatedly
selects whichever candidate is best at the current prefix, the selected
candidate can change at every arrival.  Moreover, randomized \Pivot{} offers
only an expectation at a fixed prefix and cannot provide the persistent,
pathwise control needed for low recourse.

For each sampled triple $\tau$, we instead run one deterministic, persistent
copy of Online Agreement on the evolving majority graph.  Let $A_\tau(u)$ be
its output at time $u$.  To analyze the complete history through time $t$, we
fix one optimum $F_t^*$ at the final prefix and compare every earlier output
$A_\tau(u)$ with the restriction $F_t^*|_{V_u}$.  The majority-error sets are
nested under restriction, and the exact slack identity yields
\[
  \max_{u\le t}\cost_u(A_\tau(u))
  \le
  \OPT_t+\frac{\eta+1}{2}\Delta_\tau^t(F_t^*).
\]
This is the history-stability lemma.  It is the bridge between the static
sampled-triple principle and an all-prefix online guarantee.

\paragraph{Rebuilding from monotone scores.}
For a candidate $a$, let $f_a(t)$ be its current consensus cost and let
\[
  G_a(t)=\max_{u\le t}f_a(u)
\]
be its history score.  The scores $G_a(t)$ are monotone even though a dynamic
candidate's current cost need not be.  Let $M_t=\min_aG_a(t)$.  A phase begins
with baseline $B=M_t$, follows an eligible candidate while its history score is
below $(1+\theta)B$, and discards it for the remainder of the phase once it
crosses that threshold.  A new phase starts only after $M_t$ grows by a factor
$1+\theta$.  Monotonicity implies at most one discard per candidate per phase,
and the discretization of the objective gives only
$O(\theta^{-1}\log(mn))$ phases.  Combining outer switches with the logarithmic
internal recourse of Online Agreement proves \cref{thm:online-main}.

\paragraph{Deletions through snapshot rebuilding.}
For the dynamic algorithm, we use the linear-time primitives directly.
At initialization, we select candidate input indices, candidate triples, and
an independent evaluation sample.  After every insertion or deletion, we
restrict these samples to the active set, run fresh implicit \Pivot{}
executions, and recompute every candidate's empirical cost.  With
$q=O(\delta^{-1}L)$ candidates of each type and an evaluation sample of size
$s=O(\delta^{-2}L)$, the total work is $O(qs(n_t+1))$.  Persistent input
labels can have large values; comparison sorting first compacts the sampled
labels to ranks, after which the linear-time primitives apply.

For an oblivious stream, each current instance is fixed independently of the
samples.  Its failure probability can be made at most $\zeta/(T+1)$.
Reusing the evaluation sample across different times causes no problem:
at each fixed time it remains independent of that time's candidate family,
and a union bound does not require independence between times.  The proof
never compares the current optimum with an earlier optimum.  This is why
arbitrary deletions are compatible with snapshot rebuilding although they
invalidate the monotone-history argument used for low recourse.

\paragraph{Adaptivity and the cost of fresh samples.}
If the next active instance is chosen using past outputs, it may depend on a
fixed hidden sample through those outputs.  We therefore draw all candidate
and evaluation samples anew after the adaptive adversary commits its update.
Conditioned on the past and the chosen update, the current instance is fixed
and the snapshot guarantee applies.  A conditional probability bound followed
by a union bound proves \cref{thm:adaptive-main}.  Fresh input indices can name
partitions that have never been stored, however.  The adaptive construction
therefore needs access to all current input arrays, and its small working
space is additional to that database.  This access requirement is part of the
theorem, rather than an uncharged reconstruction step.

\subsection{Further related work}
\label{subsec:related}

Consensus clustering was studied early under the name clustering
aggregation~\cite{FilkovSkiena2003,ACN2008}.  Bonizzoni, Della
Vedova, Dondi, and Jiang~\cite{BDDJ2008} established hardness and approximation
results for a constant number of input partitions.  Ailon, Charikar, and
Newman~\cite{ACN2008} developed the majority, pivot, and LP viewpoints that
underlie the strongest classical guarantees.  Cordner and
Kollios~\cite{CordnerKollios2023} developed input-label implementations and
sampling methods with approximation tradeoffs for large consensus instances.

Our local-to-global argument builds on Indyk's sublinear metric-median
sampling~\cite{Indyk1999,Indyk2001}, its development for permutation medians by
Chakraborty, Das, and Krauthgamer~\cite{CDK2021,CDK2023} and, in a unified
rank-aggregation framework, by Carmel, Das, and
Nguyen~\cite{carmel2026scalable}.  The partition-specific ingredient is the
exact overlap identity and resulting half-slack bound, together with an
output-sensitive implementation of the resulting majority graph.

For online correlation clustering, Cohen-Addad et al.~\cite{cohen2022online}
proved that a constant approximation and $O(\log n)$ worst-case per-vertex
recourse can be achieved simultaneously; they also gave a matching logarithmic
lower bound for maintaining any fixed constant approximation.  We use their
algorithm as a persistent local solver, but the sub-$2$ consensus guarantee and
the history-score rebuilding analysis are specific to the structured family of
triple-majority instances.

Dynamic unweighted correlation clustering has also been studied under vertex
and edge updates.  Cohen-Addad et al.~\cite{CLMP2024Dynamic} maintain a
constant approximation in a graph-query model; their expected amortized
polylogarithmic update-time analysis uses adversarial vertex insertions and
uniformly random vertex deletions.  Cao et al.~\cite{CaoEtAl2026Dynamic}
give an adaptive-adversary transformation with a $1.485$-approximation and
$O(\log(1/\beta))$ worst-case time per edge update, with failure probability
at most $\beta$ at any specified update.  These results concern unweighted
signed graphs.  Our dynamic construction operates directly on the label
representation of consensus inputs and accounts for sampled input access and
materializing the current consensus partition.  Its update bound follows by
recomputing the near-linear snapshot algorithm after each update; it does not
constitute a faster general algorithm for dynamic unweighted correlation
clustering.

\paragraph{Organization.}
\Cref{sec:preliminaries} introduces the partition metric, the weighted
correlation-clustering formulation, the offline access model, and the online
vertex-arrival model.  \Cref{sec:shared} proves the slack, majority, and
universal sampling lemmas used throughout.  \Cref{sec:online} presents
the sub-$2$ online algorithm, including history stability, rebuilding, and the
recourse analysis.  \Cref{sec:offline} gives the near-linear offline algorithm,
its implicit majority-graph implementation, and its approximation analysis.
\Cref{sec:dynamic} proves the fully dynamic theorems, including the
arbitrary-confidence snapshot bound, the update and space accounting, and
the adaptive extension.  We conclude in \cref{sec:conclusion}.
% =============================================================================
\section{Preliminaries}
\label{sec:preliminaries}
% =============================================================================

In this section, we introduce the common notation and the two computational
models.  All pairs are unordered and contain two distinct vertices.

\subsection{Partitions and pair-disagreement distance}

Let $W$ be a finite ground set.  A \emph{clustering} or \emph{partition} $C$
of $W$ is a family of nonempty, pairwise disjoint subsets whose union is $W$.
Its members are called clusters.  We write $u\sim_Cv$ if $u$ and $v$ lie in
the same cluster and $u\not\sim_Cv$ otherwise.  Define
\[
  E_W^+(C)=\bigl\{\{u,v\}\in\tbinom{W}{2}:u\sim_Cv\bigr\},
  \qquad
  E_W^-(C)=\tbinom{W}{2}\setminus E_W^+(C).
\]
When the ground set is clear, we omit the subscript $W$.

\begin{definition}[Partition distance]
For two partitions $C,C'$ of the same ground set $W$, their pair-disagreement
distance is
\begin{align}
  \dist_W(C,C')
  &:=\bigl|E_W^+(C)\mathbin{\triangle}E_W^+(C')\bigr| \label{eq:distance}\\
  &=\bigl|\{\{u,v\}\in\tbinom{W}{2}:
       \1[u\sim_Cv]\ne\1[u\sim_{C'}v]\}\bigr|.\notag
\end{align}
\end{definition}

Mapping $C$ to the incidence vector of $E_W^+(C)$ realizes $\dist_W$ as
Hamming distance in $\{0,1\}^{\binom{|W|}{2}}$.  Hence it is a metric.  We
will also use the Hamming triangle inequality with an arbitrary sign vector as
the intermediate point; that vector need not itself encode a transitive
partition relation.

For $U\subseteq W$, the restriction $C|_U$ is obtained by intersecting every
cluster of $C$ with $U$ and deleting empty intersections.  Restriction can only
remove disagreeing pairs, so
\begin{equation}
  \dist_U(C|_U,C'|_U)\le \dist_W(C,C').
  \label{eq:restriction-monotonicity}
\end{equation}

\subsection{Consensus clustering as a metric median}

Let $\Pp=(P_1,\ldots,P_m)$ be an indexed multiset of partitions of a common
ground set $V$, where $|V|=n$.  Repetitions are allowed.  For a partition $C$,
define its normalized consensus cost by
\begin{equation}
  \cost_{\Pp}(C)
  :=\frac1m\sum_{i=1}^{m}\dist_V(C,P_i).
  \label{eq:consensus-cost}
\end{equation}
Let
\[
  F^*\in\argmin_C\cost_{\Pp}(C),
  \qquad
  \OPT:=\cost_{\Pp}(F^*).
\]
Using the unnormalized sum multiplies every candidate and optimum cost by $m$
and therefore changes neither approximation nor competitive ratios.  We use
the normalized convention throughout so that the offline and online analyses
share exactly the same formulas.

The metric triangle inequality gives the familiar random-input bound.

\begin{observation}[Random-input benchmark]
\label{obs:random-input}
For a uniformly random index $i\in[m]$,
\[
  \E_i[\cost_{\Pp}(P_i)]
  \le 2\left(1-\frac1m\right)\OPT.
\]
\end{observation}

\begin{proof}
For every ordered pair $i\ne j$,
$\dist(P_i,P_j)\le\dist(P_i,F^*)+\dist(F^*,P_j)$.  Sum this inequality over
all ordered pairs and divide by $m^2$.  Each $\dist(P_i,F^*)$ appears $m-1$
times in each of the two symmetric terms, giving the stated factor.
\end{proof}

\subsection{Weighted correlation-clustering representation}

For a pair $e=\{u,v\}$, let
\[
  p_e:=\frac1m\bigl|\{i:u\not\sim_{P_i}v\}\bigr|.
\]
The values $p_e$ form a semimetric on $V$, because every input partition
induces a $0/1$ semimetric and averages preserve the triangle inequality.
The contribution of $e$ to \eqref{eq:consensus-cost} is $p_e$ if a solution
co-clusters the pair and $1-p_e$ if it separates the pair.  Thus consensus
clustering is metric-weighted correlation clustering on a complete graph.

More explicitly, fix any convention for the sign of a tied pair and let $H$
be the majority-signed graph of all $m$ inputs.  Put
$B_0:=\sum_{e\in\binom V2}\min\{p_e,1-p_e\}$.  A pairwise calculation gives
\begin{equation}
  \cost_{\Pp}(C)
  =B_0+
   \sum_{e\in E^+(C)\mathbin{\triangle}E^+(H)}|1-2p_e|.
  \label{eq:weighted-baseline}
\end{equation}
Indeed, following a majority sign incurs the smaller of $p_e$ and $1-p_e$;
reversing it incurs the larger and adds their absolute difference.  Ties
have zero margin and so do not depend on the chosen convention.  The
baseline $B_0$ is independent of $C$, but it matters to multiplicative
approximation guarantees.  In particular, an approximation to the total
metric-weighted objective need not approximate the number of majority
mistakes.  Our local-to-global argument will use the latter objective only
to construct candidates and then evaluate those candidates for the original
consensus cost.

We will not construct this dense graph.  Instead, our local solvers use a
constant-size sample of inputs.

\begin{definition}[Triple-majority graph]
Let $Q=\{i,j,k\}\in\binom{[m]}{3}$ be three distinct input indices.  The
\emph{majority graph} $H_Q$ is the complete signed graph on $V$ in which
$\{u,v\}$ is positive if at least two of $P_i,P_j,P_k$ co-cluster $u$ and
$v$, and negative otherwise.  We identify $H_Q$ with its set $E^+(H_Q)$ of
positive pairs.
\end{definition}

For a complete signed graph $H$ and a partition $C$, its unweighted
correlation-clustering error is
\begin{equation}
  \err_H(C):=
  |E^+(H)\cap E^-(C)|+|E^-(H)\cap E^+(C)|
  =|E^+(H)\mathbin{\triangle}E^+(C)|.
  \label{eq:cc-error}
\end{equation}
Let $\OPT_{\mathrm{CC}}(H)=\min_C\err_H(C)$.

We use two correlation-clustering primitives.  The first is the classical
randomized pivot algorithm of Ailon--Charikar--Newman~\cite{ACN2008}.

\begin{theorem}[Pivot guarantee~\cite{ACN2008}]
\label{thm:pivot-blackbox}
For every unweighted complete signed graph $H$, the partition produced by
\Pivot{} satisfies
\[
  \E[\err_H(\Pivot(H))]\le 3\OPT_{\mathrm{CC}}(H).
\]
\end{theorem}

The second primitive is deterministic and persistent under vertex arrivals.

\begin{theorem}[Online Agreement~\cite{cohen2022online}]
\label{thm:agree-blackbox}
There are absolute constants $\eta\ge1$ and $q_{\mathrm A}>0$ and a
deterministic algorithm $\Agree$ for unweighted complete signed correlation
clustering under vertex arrivals such that, at every prefix $t$,
\[
  \err_{H_t}(\Agree(H_t))
  \le\eta\OPT_{\mathrm{CC}}(H_t).
\]
Moreover, every vertex changes its persistent cluster identifier at most
$q_{\mathrm A}\lceil\log_2 n\rceil$ times over the entire execution.
\end{theorem}

This statement packages Theorems E.2 and F.1 of
Cohen-Addad et al.~\cite{cohen2022online}.  Their internal agreement parameter
is fixed to a sufficiently small absolute constant; thereafter $\eta$ and
$q_{\mathrm A}$ are fixed absolute constants.

\subsection{Offline access model and computational conventions}

An input partition $P_i$ is stored by a label array
$c^i:V\to[n]$ such that $c^i(u)=c^i(v)$ exactly when
$u\sim_{P_i}v$.  The algorithm may sample a uniform index in $[m]$ in
$O(1)$ time and then read the corresponding array in $O(n)$ time.  We count
only data actually accessed by the algorithm and report working space in
addition to the read-only input representation.  We use a word RAM with word
size sufficient to hold input indices, labels, counters, and memory addresses
in one word; in particular it is $\Omega(\log(m(n+1)))$.  For arbitrary
accuracy and confidence parameters it must also hold sample indices and
unnormalized empirical scores.  We assume a unit-cost primitive for drawing
a uniform integer in a specified word-sized range.  Consequently, an algorithm
that reads $\polylog(n)$ partitions and performs $\polylog(n)$ scans of
length $n$ has running time $\wtO(n)$ even though the full explicit input has
size $mn$.

We say an event holds \emph{with high probability} if its probability is at
least $1-n^{-c}$ for an arbitrary desired constant $c>0$, after adjusting
hidden sampling constants.  The notation $\wtO_{\eps}(\cdot)$ suppresses
polylogarithmic factors and factors depending only on $\eps$.

\subsection{Online model and recourse}

Fix an arrival ordering $v_1,\ldots,v_n$ of $V$ and let
$V_t=\{v_1,\ldots,v_t\}$.  The input consists of full partitions
$P_1,\ldots,P_m$ of $V$, but they are revealed only through their restrictions
$P_i^t:=P_i|_{V_t}$.  When $v_t$ arrives, the algorithm learns its label in
every input, equivalently its pair relation with every earlier vertex in each
input partition.  Write $\Pp^t:=(P_1^t,\ldots,P_m^t)$.

At time $t$, an online algorithm outputs a partition $C_t$ of $V_t$.  Define
\begin{equation}
  \cost_t(C):=\frac1m\sum_{i=1}^{m}\dist_{V_t}(C,P_i^t),
  \qquad
  \OPT_t:=\min_C\cost_t(C).
  \label{eq:online-cost}
\end{equation}
The optimum is monotone:
\begin{equation}
  \OPT_t\le\OPT_{t+1}.
  \label{eq:opt-monotone}
\end{equation}
Indeed, restrict an optimum at time $t+1$ to $V_t$ and use
\eqref{eq:restriction-monotonicity}.

Every output cluster carries a persistent identifier.  The first identifier
assigned to a new vertex is free; an old vertex incurs one unit of recourse
whenever its identifier changes between two consecutive prefixes.  Candidate
identifiers are namespaced, so an outer switch between distinct candidate
processes causes at most one unit of recourse to each currently present vertex.

An arrival instance is \emph{oblivious} if the full input partitions and the
arrival order are fixed independently of the algorithm's initial random
samples.  All probability statements for the online result are against such
instances.

If $m=1$, following $P_1^t$ is optimal with zero recourse.  If $m=2$, the
triangle inequality shows that $P_1^t$ is optimal: every $C$ satisfies
\[
  \dist(P_1^t,P_2^t)
  \le\dist(P_1^t,C)+\dist(C,P_2^t).
\]
Hence all the algorithms may assume $m\ge3$ after handling these cases.

\subsection{Fully dynamic model and input access}
\label{subsec:dynamic-model}

The dynamic process starts from an empty set, or from an initial instance
$\Pp^0=(P_1^0,\ldots,P_m^0)$ on $V_0$ supplied for initialization.  The
number $m\ge1$ of input indices remains fixed.  At update $t$, one of the
following operations is performed:
\begin{enumerate}[(i)]
  \item \emph{Insert a vertex $v$.}  A previously inactive vertex identifier
        is supplied with labels $(c_t^1(v),\ldots,c_t^m(v))$.  In input $i$,
        equal labels specify membership in the same cluster.  Labels of all
        previously active vertices remain unchanged.
  \item \emph{Delete a vertex $v\in V_{t-1}$.}  Remove $v$ and its label
        entries.  Every input partition is restricted to $V_{t-1}\setminus\{v\}$.
\end{enumerate}
Labels are persistent word-sized identifiers, and need not be bounded by
$n_t=|V_t|$.  Reinsertions are allowed and are treated as new insertions with
their supplied label vectors.  The theorem concerns vertex updates; it does
not silently allow an insertion to change pair relations between surviving
vertices.  Write $\Pp^t=(P_1^t,\ldots,P_m^t)$ and use the same objective
$\cost_t(C)=m^{-1}\sum_i\dist_{V_t}(C,P_i^t)$ and optimum $\OPT_t$ as in
\eqref{eq:online-cost}.  Unlike in the insertion-only model,
\eqref{eq:opt-monotone} need not hold.

In the \emph{sampled-coordinate model}, an insertion permits random access to
its $m$ labels while the update is being processed.  The algorithm pays for
the coordinates it reads and retains selected coordinates for future use.
It need not retain the entire input collection.  If the communication or
parsing interface instead forces a scan of all $m$ entries, that scan is
charged separately.  An initial nonempty instance is accessed in the same
way, one sampled partition array at a time.

In the \emph{full-current-input model}, after every update the algorithm can
read the current label array of any input index on all active vertices.
We explicitly report whether this database is external read-only input for
the clustering computation or is stored and maintained by the algorithm.
This distinction is necessary when new input indices are sampled after
earlier vertices have arrived.

An \emph{oblivious} update sequence is fixed independently of all random
choices.  An \emph{adaptive} adversary may use all previous public outputs
when choosing the next update, but must commit the new instance before the
algorithm draws that update's fresh private random choices.  In both models,
we allow arbitrary choices of the vertex to delete, subject to the stated
information available to the adversary.  Our bounds use a known capacity
$N\ge1$ and horizon $T\ge1$ and include the optional initial snapshot in
the failure budget.  Logarithms in sampling formulas are natural logarithms;
other bases affect only absolute constants.

% =============================================================================
\section{Sampled Triples and Majority Fitting}
\label{sec:shared}
% =============================================================================

In this section, we prove the common structural lemmas.  We first establish an
exact identity between triangle slack and intersections of error sets.  We then
relate the error of a triple-majority relation to total slack and prove the
universal local-to-global averaging lemma.

\subsection{The exact overlap identity}

Fix a partition $F$ and, for every input index $i$, define its error set
relative to $F$ by
\begin{equation}
  I_i(F):=E^+(P_i)\mathbin{\triangle}E^+(F).
  \label{eq:error-set}
\end{equation}
Then $|I_i(F)|=\dist(F,P_i)$.  For distinct indices $i,j$, define their
triangle-inequality slack through $F$ by
\begin{equation}
  \Delta_{ij}(F)
  :=\dist(F,P_i)+\dist(F,P_j)-\dist(P_i,P_j)\ge0.
  \label{eq:pair-slack}
\end{equation}
For a set $Q$ of indices, let
\begin{equation}
  \Delta_Q(F):=\sum_{\{i,j\}\in\binom{Q}{2}}\Delta_{ij}(F).
  \label{eq:total-slack}
\end{equation}

\begin{lemma}[Exact overlap identity]
\label{lem:exact-overlap}
For every partition $F$ and all distinct indices $i,j\in[m]$,
\[
  \Delta_{ij}(F)=2|I_i(F)\cap I_j(F)|.
\]
\end{lemma}

\begin{proof}
Symmetric difference is associative and self-inverse, so
\[
  E^+(P_i)\mathbin{\triangle}E^+(P_j)
  =I_i(F)\mathbin{\triangle}I_j(F).
\]
Therefore
\[
  \dist(P_i,P_j)
  =|I_i(F)|+|I_j(F)|-2|I_i(F)\cap I_j(F)|.
\]
Substituting this equality into \eqref{eq:pair-slack} proves the claim.
\end{proof}

Fix a triple $Q=\{i,j,k\}$.  Let $B_Q(F)$ be the pairs on which the majority
sign of $P_i,P_j,P_k$ disagrees with $F$.

\begin{lemma}[Majority error is at most half the slack]
\label{lem:majority-half-slack}
For every partition $F$ and every triple $Q=\{i,j,k\}$,
\begin{align}
  B_Q(F)
  &=(I_i(F)\cap I_j(F))
    \cup(I_i(F)\cap I_k(F))
    \cup(I_j(F)\cap I_k(F)),
    \label{eq:majority-error-union}\\
  |B_Q(F)|&\le\frac12\Delta_Q(F).
    \label{eq:majority-half}
\end{align}
In particular,
$\err_{H_Q}(F)=|B_Q(F)|$ and
$\OPT_{\mathrm{CC}}(H_Q)\le\Delta_Q(F)/2$.
\end{lemma}

\begin{proof}
For a fixed pair of vertices, encode co-clustering by $1$ and separation by
$0$.  The majority bit of three inputs differs from the bit of $F$ exactly
when at least two input bits differ from it.  This proves
\eqref{eq:majority-error-union}.  A union bound followed by
\cref{lem:exact-overlap} gives
\[
  |B_Q(F)|
  \le\sum_{\{a,b\}\in\binom{Q}{2}}|I_a(F)\cap I_b(F)|
  =\frac12\Delta_Q(F).
\]
The last assertions follow from the definition of correlation-clustering
error and the feasibility of $F$ for $H_Q$.
\end{proof}

The next consequence will be instantiated with either \Pivot{} or \Agree{}.

\begin{lemma}[Majority bridge]
\label{lem:majority-bridge}
Let $A$ be any partition satisfying
$\err_{H_Q}(A)\le\gamma\OPT_{\mathrm{CC}}(H_Q)$.  Then, for every reference
partition $F$,
\[
  \dist(F,A)
  \le\frac{\gamma+1}{2}\Delta_Q(F).
\]
\end{lemma}

\begin{proof}
The Hamming triangle inequality through the majority sign vector gives
\[
  \dist(F,A)
  \le\err_{H_Q}(F)+\err_{H_Q}(A).
\]
By \cref{lem:majority-half-slack},
$\err_{H_Q}(F)\le\Delta_Q(F)/2$ and
$\OPT_{\mathrm{CC}}(H_Q)\le\Delta_Q(F)/2$.  Substitute the assumed bound on
$A$.
\end{proof}

\subsection{A universal averaging lemma}

The following lemma is a self-contained form of the local-to-global sampling
principle used for metric medians~\cite{carmel2026scalable}.  It applies to
arbitrary metric spaces; no partition structure is used.

\begin{lemma}[Universal slack averaging]
\label{lem:universal-averaging}
Let $P_1,\ldots,P_m$ be points in a metric space, let $F^*$ be an optimal
$1$-median, and let
$\cost_{\Pp}(x):=m^{-1}\sum_j\dist(x,P_j)$ and
$\OPT:=\cost_{\Pp}(F^*)$.  Fix $\alpha\in(0,1)$ and
$\rho\in[0,1-\alpha]$.  Suppose at most $\rho m$ inputs satisfy
\[
  \cost_{\Pp}(P_i)\le(2-\alpha)\OPT.
\]
If $Q$ is a uniformly random $r$-element subset of the indexed inputs, where
$2\le r\le m$, then
\[
  \E_Q[\Delta_Q(F^*)]
  \le\binom r2(\alpha+2\rho)\OPT.
\]
\end{lemma}

\begin{proof}
Let $G$ be the set of inputs with cost at most $(2-\alpha)\OPT$.  Since
$|G|\le\rho m$,
\begin{equation}
  \sum_{i=1}^{m}\cost_{\Pp}(P_i)
  \ge(1-\rho)m(2-\alpha)\OPT.
  \label{eq:sum-input-lb}
\end{equation}
On the other hand,
\[
  \sum_i\cost_{\Pp}(P_i)
  =\frac1m\sum_{i\ne j}\dist(P_i,P_j).
\]
Writing each pairwise distance in terms of its slack through $F^*$ gives
\begin{align*}
  \sum_{i\ne j}\Delta_{ij}(F^*)
  &=2(m-1)\sum_i\dist(F^*,P_i)
    -\sum_{i\ne j}\dist(P_i,P_j)\\
  &=2m(m-1)\OPT
    -m\sum_i\cost_{\Pp}(P_i)\\
  &\le\bigl(m^2(\alpha+2\rho-\alpha\rho)-2m\bigr)\OPT\\
  &\le\bigl(m^2(\alpha+2\rho)-2m\bigr)\OPT,
\end{align*}
where the first inequality uses \eqref{eq:sum-input-lb}.  Hence
\begin{equation}
  \sum_{1\le i<j\le m}\Delta_{ij}(F^*)
  \le\frac12\bigl(m(\alpha+2\rho)-2\bigr)m\OPT.
  \label{eq:total-slack-bound}
\end{equation}
Each fixed pair belongs to a uniform $r$-subset with probability
\[
  \frac{\binom{m-2}{r-2}}{\binom mr}
  =\frac{r(r-1)}{m(m-1)}.
\]
Multiplying \eqref{eq:total-slack-bound} by this probability gives
\[
  \E_Q[\Delta_Q(F^*)]
  \le\binom r2
     \frac{m(\alpha+2\rho)-2}{m-1}\OPT.
\]
Finally, $\rho\le1-\alpha$ implies $\alpha+2\rho\le2$, and therefore
the fraction is at most $\alpha+2\rho$.
\end{proof}

For triples, Markov's inequality turns the expectation into a useful sampling
probability.

\begin{corollary}[Finding a low-slack triple]
\label{cor:low-slack-triple}
Under the hypotheses of \cref{lem:universal-averaging} with $r=3$ and
$\rho>0$, a uniformly random triple satisfies
\[
  \Delta_Q(F^*)\le3(\alpha+3\rho)\OPT
\]
with probability at least
\[
  \frac{\rho}{\alpha+3\rho}.
\]
\end{corollary}

\begin{proof}
The expectation is at most $3(\alpha+2\rho)\OPT$.  Markov's inequality bounds
the probability of exceeding $3(\alpha+3\rho)\OPT$ by
$(\alpha+2\rho)/(\alpha+3\rho)$.
\end{proof}

% =============================================================================
\section{Sub-Two Online Consensus Clustering}
\label{sec:online}
% =============================================================================

In this section, we prove \cref{thm:online-main}.  We first define a fixed
random family of persistent candidate processes.  We then prove that, with
high probability, at every terminal prefix at least one candidate has been
good throughout its complete history.  Finally, a rebuilding rule tracks these
history-good candidates with only polylogarithmically many switches.

\subsection{Candidate processes}
\label{subsec:online-candidates}

Fix the error tolerance $\eps$ and failure probability $\zeta$ from
\cref{thm:online-main}, and define
\begin{equation}
  \alpha_*:=\frac{2}{3\eta+5},
  \qquad
  \rho:=\frac{\eps}{9(\eta+1)}.
  \label{eq:online-parameters-one}
\end{equation}
The assumption $\eps<\alpha_*$ ensures
$0<\rho<1-\alpha_*$.  Set
\begin{align}
  L_{\mathrm{in}}
  &:=\left\lceil\frac{1}{\rho}
       \ln\frac{2n}{\zeta}\right\rceil,
       \label{eq:num-input-candidates}\\
  L_{\mathrm{tri}}
  &:=\left\lceil
       \frac{\alpha_*+3\rho}{\rho}
       \ln\frac{2n}{\zeta}\right\rceil.
       \label{eq:num-triple-candidates}
\end{align}
Before the first arrival, sample $L_{\mathrm{in}}$ independent uniform input
indices and $L_{\mathrm{tri}}$ independent uniform three-element subsets of
$[m]$.  Repetitions across different samples are harmless.  Let $\A$ denote
the resulting family of candidate processes and let
\begin{equation}
  K:=|\A|\le L_{\mathrm{in}}+L_{\mathrm{tri}}
  =O_\eta\!\left(\eps^{-1}\log\frac{2n}{\zeta}\right).
  \label{eq:candidate-count}
\end{equation}

There are two candidate types.

\paragraph{Input candidates.}
For a sampled index $i$, define
\[
  A_i(t):=P_i^t.
\]
The cluster labels of this process are inherited from the fixed full
partition $P_i$.  Hence it has zero internal recourse.

\paragraph{Triple candidates.}
For a sampled triple $Q=\{i,j,k\}$, let $H_Q(t)$ be the majority signed graph
of $P_i^t,P_j^t,P_k^t$ on $V_t$.  These graphs form a valid vertex-arrival
sequence: $H_Q(t)$ is the subgraph of $H_Q(t+1)$ induced by $V_t$.  Run one
persistent deterministic copy of Online Agreement and define
\[
  A_Q(t):=\Agree(H_Q(t)).
\]
By \cref{thm:agree-blackbox}, this candidate is $\eta$-approximate for the
auxiliary signed instance at every prefix and causes at most
$q_{\mathrm A}\lceil\log_2 n\rceil$ internal identifier changes to any
vertex over the full execution.

For every candidate $a\in\A$, define its current score and history score by
\begin{equation}
  f_a(t):=\cost_t(A_a(t)),
  \qquad
  G_a(t):=\max_{1\le u\le t}f_a(u).
  \label{eq:history-score}
\end{equation}
The quantity $G_a(t)$ is nondecreasing.  For an input candidate $i$, each
pairwise distance $\dist_{V_t}(P_i^t,P_h^t)$ is nondecreasing with $t$;
therefore
\begin{equation}
  G_i(t)=f_i(t).
  \label{eq:input-current-history}
\end{equation}
Finally, let
\begin{equation}
  M_t:=\min_{a\in\A}G_a(t).
  \label{eq:min-history-score}
\end{equation}
Since it is the pointwise minimum of finitely many nondecreasing functions,
$M_t$ is also nondecreasing.

\begin{paperalgorithm}[htbp]
\centering
\fbox{\begin{minipage}{0.93\textwidth}
\textbf{Sampled online candidate family.}

\begin{enumerate}[(1)]
  \item Before the stream, sample the input indices and distinct-index triples
        in \eqref{eq:num-input-candidates}--\eqref{eq:num-triple-candidates}.
  \item For every sampled input, maintain its restriction to the current
        prefix.  For every sampled triple, maintain a persistent Online
        Agreement execution on its evolving majority graph.
  \item At every time $t$, update the current scores $f_a(t)$, history scores
        $G_a(t)$, and their minimum $M_t$.
  \item Use the rebuilding rule in \cref{subsec:rebuilding} to select one
        candidate process; output that candidate's current partition.
\end{enumerate}
\end{minipage}}
\caption{The sampled candidate processes used by the online algorithm.}
\label{alg:online-candidates}
\end{paperalgorithm}

\subsection{History stability of a triple candidate}
\label{subsec:history-stability}

The crucial feature of Online Agreement is not merely that it is approximate
at the current time.  Its deterministic all-prefix guarantee lets us compare
the entire trajectory of a triple candidate with restrictions of one fixed
terminal-prefix optimum.

For a partition $F$ of $V_t$, define the prefix-restricted quantities
\begin{align*}
  I_i^t(F)
  &:=E_{V_t}^+(P_i^t)\mathbin{\triangle}E_{V_t}^+(F),\\
  \Delta_{ij}^t(F)
  &:=\dist_{V_t}(F,P_i^t)+\dist_{V_t}(F,P_j^t)
     -\dist_{V_t}(P_i^t,P_j^t),\\
  \Delta_Q^t(F)
  &:=\sum_{\{i,j\}\in\binom Q2}\Delta_{ij}^t(F),
\end{align*}
and let $B_Q^t(F)$ be the pairs on which the majority of the restricted
inputs indexed by $Q$ disagrees with $F$.

Fix a terminal prefix $t$ and an optimum $F_t^*$ of $\Pp^t$.  For a triple
$Q$ and an earlier time $u\le t$, let
\begin{align*}
  B_{Q,t,t}&:=B_Q^t(F_t^*),\\
  B_{Q,t,u}&:=B_{Q,t,t}\cap\binom{V_u}{2}.
\end{align*}
Here the first set is formed on $V_t$ and the second restricts it to earlier
pairs.  Both use the same comparator $F_t^*$; we never replace it by an
unrelated optimum $F_u^*$.

Because majority relations commute with restriction, $B_{Q,t,u}$ is exactly
the set of pairs on which $H_Q(u)$ disagrees with $F_t^*|_{V_u}$.  Therefore
\begin{equation}
  \err_{H_Q(u)}(F_t^*|_{V_u})=|B_{Q,t,u}|.
  \label{eq:restricted-comparator-error}
\end{equation}

\begin{lemma}[History bound for one triple]
\label{lem:triple-history-bound}
For every terminal prefix $t$ and every triple $Q$,
\[
  G_Q(t)
  \le
  \OPT_t+\frac{\eta+1}{2}\Delta_Q^t(F_t^*),
\]
where $\Delta_Q^t$ is computed on the restricted inputs
$P_i^t$, $i\in Q$.
\end{lemma}

\begin{proof}
Fix $u\le t$.  Online Agreement and
\eqref{eq:restricted-comparator-error} give
\begin{align*}
  \err_{H_Q(u)}(A_Q(u))
  &\le\eta\OPT_{\mathrm{CC}}(H_Q(u))\\
  &\le\eta\err_{H_Q(u)}(F_t^*|_{V_u})
   =\eta|B_{Q,t,u}|.
\end{align*}
The Hamming triangle inequality through the majority sign vector yields
\begin{equation}
  \dist_{V_u}(A_Q(u),F_t^*|_{V_u})
  \le(\eta+1)|B_{Q,t,u}|.
  \label{eq:online-distance-to-terminal-opt}
\end{equation}
Now compare the current consensus cost with the optimum at the later time
$t$.  By the partition triangle inequality and restriction monotonicity,
\begin{align*}
  f_Q(u)
  &=\frac1m\sum_{h=1}^{m}\dist_{V_u}(A_Q(u),P_h^u)\\
  &\le\frac1m\sum_{h=1}^{m}
       \dist_{V_u}(F_t^*|_{V_u},P_h^u)
       +\dist_{V_u}(A_Q(u),F_t^*|_{V_u})\\
  &\le\OPT_t+(\eta+1)|B_{Q,t,u}|\\
  &\le\OPT_t+(\eta+1)|B_{Q,t,t}|\\
  &\le\OPT_t+\frac{\eta+1}{2}\Delta_Q^t(F_t^*),
\end{align*}
where the last line uses \cref{lem:majority-half-slack} on the terminal
prefix.  The bound holds for every $u\le t$; taking their maximum proves the
lemma.
\end{proof}

\subsection{A history-good candidate at every prefix}
\label{subsec:simultaneous-candidates}

We now apply the metric averaging lemma separately to each fixed prefix.  The
oblivious-adversary assumption is used here: once the full instance is fixed,
the candidate samples are independent of the optimum and costs at every
prefix.

\begin{lemma}[Fixed-prefix candidate dichotomy]
\label{lem:fixed-prefix-dichotomy}
Fix a time $t$.  With probability at least $1-\zeta/n$ over the initial
candidate sampling, the family $\A$ contains a candidate $a$ satisfying
\begin{equation}
  G_a(t)
  \le
  \left(2-\alpha_*+\frac{\eps}{2}\right)\OPT_t.
  \label{eq:fixed-prefix-history-good}
\end{equation}
\end{lemma}

\begin{proof}
Let
\[
  \mathcal{G}_t
  :=\bigl\{i\in[m]:
       \cost_t(P_i^t)\le(2-\alpha_*)\OPT_t\bigr\}.
\]
We distinguish two cases.

\emph{Many good inputs.}
If $|\mathcal{G}_t|\ge\rho m$, then one uniform input sample belongs to
$\mathcal{G}_t$ with probability at least $\rho$.  By
\eqref{eq:num-input-candidates}, the probability that all input samples miss
$\mathcal{G}_t$ is at most
\[
  (1-\rho)^{L_{\mathrm{in}}}
  \le e^{-\rho L_{\mathrm{in}}}
  \le\frac{\zeta}{2n}.
\]
For a sampled $i\in\mathcal{G}_t$,
\eqref{eq:input-current-history} gives
\[
  G_i(t)=\cost_t(P_i^t)
  \le(2-\alpha_*)\OPT_t,
\]
which is stronger than \eqref{eq:fixed-prefix-history-good}.

\emph{Few good inputs.}
Suppose $|\mathcal{G}_t|<\rho m$.  Apply
\cref{cor:low-slack-triple} in the partition metric at prefix $t$, with
$\alpha=\alpha_*$.  A uniform triple satisfies
\begin{equation}
  \Delta_Q^t(F_t^*)
  \le3(\alpha_*+3\rho)\OPT_t
  \label{eq:online-low-slack-event}
\end{equation}
with probability at least $\rho/(\alpha_*+3\rho)$.  The probability that all
$L_{\mathrm{tri}}$ independent triple samples fail is at most
\[
  \left(1-\frac{\rho}{\alpha_*+3\rho}\right)^{L_{\mathrm{tri}}}
  \le\frac{\zeta}{2n}.
\]
For a triple satisfying \eqref{eq:online-low-slack-event},
\cref{lem:triple-history-bound} gives
\begin{align*}
  G_Q(t)
  &\le\left(
    1+\frac{3(\eta+1)}2(\alpha_*+3\rho)
  \right)\OPT_t\\
  &=\left(
    2-\alpha_*+\frac{9(\eta+1)}2\rho
  \right)\OPT_t\\
  &=\left(2-\alpha_*+\frac{\eps}{2}\right)\OPT_t.
\end{align*}
The second line uses
\begin{equation}
  \alpha_*=\frac{1}{1+3(\eta+1)/2}
  =\frac{2}{3\eta+5},
  \label{eq:online-balance}
\end{equation}
and the last line uses the definition of $\rho$ in
\eqref{eq:online-parameters-one}.
\end{proof}

A union bound now makes the result simultaneous.

\begin{lemma}[Simultaneous history event]
\label{lem:simultaneous-history-event}
With probability at least $1-\zeta$, simultaneously for every
$t\in[n]$,
\begin{equation}
  M_t
  \le R\OPT_t,
  \qquad
  R:=2-\alpha_*+\frac{\eps}{2}.
  \label{eq:simultaneous-M}
\end{equation}
\end{lemma}

\begin{proof}
Apply \cref{lem:fixed-prefix-dichotomy} to each of the $n$ fixed prefixes and
take a union bound.  Whenever the lemma succeeds at $t$, the witness candidate
has $G_a(t)\le R\OPT_t$, so the minimum $M_t$ does as well.
\end{proof}

No independence between different prefixes is asserted or needed.  The same
fixed candidate family is used at every prefix.

\subsection{The rebuilding algorithm}
\label{subsec:rebuilding}

Set the hysteresis parameter to
\begin{equation}
  \theta:=\frac{\eps}{4}.
  \label{eq:hysteresis}
\end{equation}
The output follows one current candidate.  Rebuilding is driven by the
monotone scores $G_a(t)$ and $M_t$ rather than current costs.

\begin{paperalgorithm}[htbp]
\centering
\fbox{\begin{minipage}{0.93\textwidth}
\textbf{History-score rebuilding.}

\begin{enumerate}[(1)]
  \item If $M_t=0$, retain the current candidate if its history score is zero.
        Otherwise, select a candidate of history score zero (also doing so
        if no candidate has yet been selected).  Switch during this stage
        only when the selected candidate's score becomes positive.
  \item When $M_t$ first becomes positive, start a positive phase.  Set its
        baseline $B:=M_t$, clear its discarded-candidate set, and select a
        minimizer of $G_a(t)$.
  \item A positive phase continues while $M_t<(1+\theta)B$.  Within it,
        follow the current candidate $a$ while
        $G_a(t)<(1+\theta)B$.
  \item At the first time $M_t\ge(1+\theta)B$, end the phase and begin a new
        one with baseline $M_t$ and a new minimizer.  This phase-boundary test
        takes precedence over the next step.
  \item Otherwise, if the current candidate reaches the threshold, discard it
        for the remainder of the phase and switch to any undiscarded candidate
        whose history score is below $(1+\theta)B$.
\end{enumerate}
\end{minipage}}
\caption{The monotone rebuilding rule.}
\label{alg:online-rebuilding}
\end{paperalgorithm}

The switch in Step~(5) is always possible.  While a phase continues, a
candidate attaining $M_t$ has score below $(1+\theta)B$.  A candidate discarded
earlier has score at least the threshold forever, because history scores are
nondecreasing.  Hence the current minimizer has not been discarded.

\subsection{Competitive analysis}
\label{subsec:online-competitive}

\begin{lemma}[Cost within a phase]
\label{lem:online-phase-cost}
Assume the simultaneous history event of
\cref{lem:simultaneous-history-event}.  At every time $t$, the rebuilding
algorithm outputs a partition satisfying
\[
  \cost_t(C_t)\le(1+\theta)R\OPT_t.
\]
When $M_t=0$, its output cost is zero.
\end{lemma}

\begin{proof}
If $M_t=0$, the algorithm follows a candidate with history score zero, whose
current cost is also zero.

Consider a positive phase that starts at time $t_0$ with baseline
$B=M_{t_0}$.  Throughout the phase, the selected candidate $a$ satisfies
\[
  f_a(t)\le G_a(t)<(1+\theta)B.
\]
By the simultaneous event and monotonicity of the optimum,
\[
  B=M_{t_0}\le R\OPT_{t_0}\le R\OPT_t.
\]
Substitution proves the claim.  At a phase boundary, the new minimizer is
selected before the output is declared, so the same argument applies.
\end{proof}

Since $R=2-\alpha_*+\eps/2<2$ and $\theta=\eps/4$,
\begin{equation}
  (1+\theta)R
  =R+\theta R
  <2-\alpha_*+\frac{\eps}{2}+\frac{\eps}{2}
  =2-\frac{2}{3\eta+5}+\eps.
  \label{eq:online-final-ratio}
\end{equation}
Together with \cref{lem:simultaneous-history-event}, this proves the
competitive-ratio part of \cref{thm:online-main}.

\subsection{Number of rebuilds and per-vertex recourse}
\label{subsec:online-recourse}

The case $n=1$ has zero cost and zero recourse, so assume $n\ge2$.  Every
normalized candidate cost belongs to
\[
  \left\{0,\frac1m,\frac2m,\ldots,\binom n2\right\}.
\]
The same is true of every history score and of $M_t$.  During the zero stage,
a candidate that becomes positive can never again have history score zero.
Hence at most $K$ candidate switches occur before $M_t$ becomes positive.

Let $B_1,\ldots,B_s$ be the baselines of the positive phases.  The phase rule
ensures $B_{h+1}\ge(1+\theta)B_h$.  Every positive baseline lies between
$1/m$ and $\binom n2$, so
\begin{align}
  s
  &\le 1+\left\lceil
      \log_{1+\theta}\left(m\binom n2\right)
    \right\rceil
   =O\!\left(\eps^{-1}\log(2mn)\right).
  \label{eq:number-online-phases}
\end{align}
Within one phase, a discarded candidate remains discarded.  Thus, there are
at most $K$ switches per phase, up to an inessential boundary choice.  The
total number of outer switches is
\begin{equation}
  O\!\left(K\left[1+\eps^{-1}\log(2mn)\right]\right).
  \label{eq:outer-switch-count}
\end{equation}

Fix a vertex $v$.  Each outer switch changes its namespaced identifier at
most once.  An input candidate causes no internal recourse.  For each triple
candidate, \cref{thm:agree-blackbox} bounds the internal recourse of $v$ over
the full stream by $q_{\mathrm A}\lceil\log_2 n\rceil$.  Re-entering a
candidate is already charged as an outer switch.  Therefore
\begin{align}
  \operatorname{rec}(v)
  &\le O\!\left(
      K\left[
        1+\eps^{-1}\log(2mn)+q_{\mathrm A}\log(2n)
      \right]
    \right)\\
  &=O_{\eta,q_{\mathrm A}}\!\left(
      \eps^{-2}
      \log\frac{2n}{\zeta}
      \log\frac{2mn}{\zeta}
    \right),
  \label{eq:online-recourse-final}
\end{align}
where \eqref{eq:candidate-count} and $\eps<1$ are used in the second line.
This analysis is deterministic once the initial sample family is fixed.
Combining \eqref{eq:online-final-ratio} and
\eqref{eq:online-recourse-final} completes the proof of
\cref{thm:online-main}.

\subsection{Scope of the online result}

The proof requires an oblivious instance because the same initial samples are
analyzed by a union bound over fixed prefixes.  It also assumes that $n$ is
known when the candidate family is sampled.  We do not claim the usual
unknown-horizon doubling extension here, because a restart would have to
initialize persistent candidate histories and identifiers carefully.

The algorithm maintains exact scores against all $m$ inputs.  This is
polynomial time, but the theorem is a statement about competitive ratio and
recourse, not update time or memory.  The proof uses a persistent local solver
with an all-prefix approximation guarantee and bounded per-vertex recourse.
Deterministic Online Agreement supplies this interface.  Independently
rerunning a randomized offline algorithm does not supply either guarantee.

% =============================================================================
\section{Near-Linear-Time Offline Consensus Clustering}
\label{sec:offline}
% =============================================================================

In this section, we prove \cref{thm:offline-main}.  The algorithm instantiates
the local-to-global framework of \cref{lem:universal-averaging} with triples.
Its local solver runs \Pivot{} on a triple-majority graph.  The main
implementation challenge is that this graph can contain $\Theta(n^2)$ positive
pairs; we give an implicit representation on which the complete \Pivot{}
execution takes $O(n)$ time.

\subsection{The global sampling algorithm}
\label{subsec:offline-global}

If $n\le1$, the unique partition has zero cost.  If $m\le2$, return $P_1$, which
is optimal by the argument in \cref{sec:preliminaries}.  Henceforth assume
$n\ge2$ and $m\ge3$.

We first state the algorithm with an internal accuracy parameter
$\delta\in(0,1/14]$.  The final choice of $\delta$ in terms of the theorem's
$\eps$ is made at the end of the analysis.  Let
\begin{equation}
  q_{\mathrm{in}}=\Theta(\delta^{-1}\log n),
  \qquad
  q_{\mathrm{loc}}=\Theta(\delta^{-1}\log n),
  \qquad
  s=\Theta\!\left(\delta^{-2}\log\frac{n}{\delta}\right),
  \label{eq:offline-sample-sizes}
\end{equation}
where the constants can be enlarged to obtain any desired polynomially small
failure probability.

The algorithm constructs a candidate family $\Y$ in two ways.  First, it
samples $q_{\mathrm{in}}$ uniform input indices and adds the corresponding
partitions.  Second, it performs $q_{\mathrm{loc}}$ independent local trials.
In each trial, it samples a uniform three-element subset
$Q\in\binom{[m]}3$, runs \Pivot{} on the majority graph $H_Q$, and adds the
output to $\Y$.  Finally, it draws a fresh independent multiset $S$ of $s$
uniform input indices and returns the candidate minimizing empirical cost on
$S$.

\begin{paperalgorithm}[htbp]
\centering
\fbox{\begin{minipage}{0.93\textwidth}
\textbf{Offline sampled-triple consensus algorithm.}

\begin{enumerate}[(1)]
  \item Initialize $\Y$ with $q_{\mathrm{in}}$ independently sampled input
        partitions.
  \item Repeat $q_{\mathrm{loc}}$ times independently:
        sample $Q\in\binom{[m]}3$ uniformly; build the implicit majority
        oracle for $H_Q$; run \Pivot{} on this oracle; add its output to
        $\Y$.
  \item Independently sample an evaluation multiset $S$ of $s$ input
        partitions.
  \item Return
        \[
          \widehat C\in\argmin_{C\in\Y}
          \frac1{|S|}\sum_{P\in S}\dist(C,P).
        \]
\end{enumerate}
\end{minipage}}
\caption{The near-linear-time offline algorithm.  The evaluation sample is
fresh and independent of the random candidate family.}
\label{alg:offline-global}
\end{paperalgorithm}

The remainder of the section supplies the two linear-time primitives and then
proves the approximation guarantee.

\subsection{Computing partition distance in linear time}
\label{subsec:linear-distance}

Let $A$ and $B$ be two partitions of $V$.  Denote their clusters by
$A_1,A_2,\ldots$ and $B_1,B_2,\ldots$.  Their meet $A\wedge B$ consists of
the nonempty intersections $A_i\cap B_j$.  Since
\[
  E^+(A)\cap E^+(B)=E^+(A\wedge B),
\]
the identity $|X\triangle Y|=|X|+|Y|-2|X\cap Y|$ gives
\begin{equation}
  \dist(A,B)
  =\sum_i\binom{|A_i|}{2}
   +\sum_j\binom{|B_j|}{2}
   -2\sum_{i,j}\binom{|A_i\cap B_j|}{2}.
  \label{eq:contingency-distance}
\end{equation}

\begin{lemma}[Linear-time distance evaluation]
\label{lem:linear-distance}
Given the label arrays of two partitions $A$ and $B$, the distance
$\dist(A,B)$ can be computed in $O(n)$ time and $O(n)$ space on a word RAM.
\end{lemma}

\begin{proof}
The sizes $|A_i|$ and $|B_j|$ are counted by two scans using arrays indexed by
cluster labels.  To count the nonempty intersections, radix-sort the vertices
by the pair key $(c^A(v),c^B(v))$.  Two stable counting-sort passes suffice
because both labels lie in $[n]$.  Equal pair keys form contiguous runs, and a
scan of these runs gives every $|A_i\cap B_j|$.  Substituting the counts into
\eqref{eq:contingency-distance} completes the computation in linear time and
space.
\end{proof}

Consequently, evaluating one candidate on $s$ sampled inputs costs $O(sn)$,
without examining the unsampled inputs.

\subsection{An implicit majority graph for three partitions}
\label{subsec:implicit-majority}

Fix a local sample $Q=\{i_1,i_2,i_3\}\in\binom{[m]}3$, and write
$W_r:=P_{i_r}$ for $r\in\{1,2,3\}$.  Let $c_r(v)\in[n]$ denote the
cluster label of $v$ in $W_r$.  Thus every vertex has a
label tuple
\[
  c(v)=(c_1(v),c_2(v),c_3(v)).
\]
Two vertices are positive in $H_Q$ precisely when their tuples agree in at
least two coordinates.

For $1\le a<b\le3$, define the pair-label bucket
\begin{equation}
  B_{ab}(v)
  :=\{u\in V:c_a(u)=c_a(v)\text{ and }c_b(u)=c_b(v)\}.
  \label{eq:pair-bucket}
\end{equation}
There are three bucket families, indexed by $(1,2)$, $(1,3)$, and $(2,3)$.

\begin{observation}[Neighborhood as three buckets]
\label{obs:three-buckets}
For every vertex $v$, its closed positive neighborhood in $H_Q$ is
\begin{equation}
  N_{H_Q}[v]=B_{12}(v)\cup B_{13}(v)\cup B_{23}(v).
  \label{eq:neighborhood-union}
\end{equation}
Every member of this union occurs in at most three of the listed buckets.
\end{observation}

\begin{proof}
A vertex $u$ belongs to the union exactly when its tuple agrees with $c(v)$
on one of the three coordinate pairs, equivalently on at least two
coordinates.  This is precisely the majority-positive condition.  There are
only three buckets, proving the multiplicity bound.
\end{proof}

We make this representation dynamic under deletions.  For each pair type
$ab$, group vertices by the key $(c_a(v),c_b(v))$ and store every nonempty
bucket as a doubly linked list of currently active vertices.  Every vertex
keeps one back-pointer to its node in each of its three lists, and every list
node stores a pointer to its bucket header so that the full bucket can be
enumerated from any member.  The buckets can be built in $O(n)$ time by
two-pass radix sorting for each of the three pair types.

We also store the active vertices in a swap array $A[0],\ldots,A[k-1]$ with
an inverse-position array.  A uniform active pivot is obtained by choosing a
uniform index in $\{0,\ldots,k-1\}$.  To delete a vertex, swap it with the last
active element and decrement $k$.  This operation, as well as splicing the
vertex out of its three bucket lists through its back-pointers, costs $O(1)$.

At a pivot step, enumerate the three active buckets of the pivot, use a
timestamp array to remove duplicates, output their union as one cluster, and
delete every output vertex from all data structures.  The following framed
description makes all operations explicit.

\begin{paperalgorithm}[htbp]
\centering
\fbox{\begin{minipage}{0.93\textwidth}
\textbf{Implicit \Pivot{} on a triple-majority graph.}

\begin{enumerate}[(1)]
  \item \textbf{Build.}  From the three label arrays, construct the active
        doubly linked buckets $B_{12},B_{13},B_{23}$, their vertex back-pointers,
        the swap array, its inverse positions, and a zeroed timestamp array.
  \item While active vertices remain, choose a uniform active pivot $p$ from
        the swap array and increment the current timestamp.
  \item Enumerate $B_{12}(p)$, $B_{13}(p)$, and $B_{23}(p)$.  Add a vertex to
        the new cluster the first time it is encountered at the current
        timestamp.
  \item Output the resulting union as the next cluster.
        Assign all its vertices a common fresh label in the output label array.
  \item For every vertex in that cluster, splice it out of its three bucket
        lists and delete it from the swap array.
\end{enumerate}
\end{minipage}}
\caption{An output-sensitive simulation of \Pivot{} that never materializes
the majority graph.}
\label{alg:implicit-pivot}
\end{paperalgorithm}

\begin{lemma}[Linear-time implicit Pivot]
\label{lem:implicit-pivot-linear}
Under the same random pivot choices, the procedure in
\cref{alg:implicit-pivot} outputs exactly the partition produced by ordinary
\Pivot{} on $H_Q$.  Its total running time and space are $O(n)$.
\end{lemma}

\begin{proof}
The back-pointer deletions maintain the invariant that every bucket contains
exactly its active vertices.  Hence, by \cref{obs:three-buckets}, the three
lists enumerated at pivot $p$ have union equal to the current closed positive
neighborhood of $p$.  Timestamps ensure each vertex is output once in the
step.  The swap-array invariant makes $p$ uniform among all active vertices.
Thus the procedure simulates ordinary \Pivot{} exactly.

Building the three bucket families costs $O(n)$.  Consider a pivot step that
outputs cluster $C$.  Every enumerated list member lies in $C$ and appears in
at most three lists, so the enumeration costs at most $O(3|C|)=O(|C|)$.
Deleting the vertices costs another $O(|C|)$.  The output clusters partition
$V$; assigning their output labels costs another $O(n)$ in total.  Therefore
\[
  \sum_{\text{output clusters }C}O(|C|)=O(n).
\]
The lists, three back-pointers per vertex, swap array, inverse positions, and
timestamps use $O(n)$ words.
\end{proof}

The proof establishes exactly the interface used by \Pivot{}: uniform active
pivots and output-sensitive enumeration of their closed neighborhoods.  We do
not claim that this data structure supports every possible graph query in
constant time.

\subsection{The local approximation property}
\label{subsec:offline-local-property}

Fix an optimum $F^*$ for the full consensus instance.  Let $Y_Q$ be the
random output of \Pivot{} on $H_Q$, with the sample $Q$ fixed.

\begin{lemma}[Local property with constant two]
\label{lem:offline-local-C2}
For every triple $Q$,
\[
  \E_{\mathrm{piv}}[\dist(F^*,Y_Q)\mid Q]
  \le 2\Delta_Q(F^*).
\]
\end{lemma}

\begin{proof}
The Hamming triangle inequality through the majority sign vector gives
\begin{equation}
  \dist(F^*,Y_Q)
  \le\err_{H_Q}(F^*)+\err_{H_Q}(Y_Q).
  \label{eq:offline-majority-bridge}
\end{equation}
By the pivot guarantee and the feasibility of $F^*$,
\[
  \E_{\mathrm{piv}}[\err_{H_Q}(Y_Q)\mid Q]
  \le3\OPT_{\mathrm{CC}}(H_Q)
  \le3\err_{H_Q}(F^*).
\]
Taking expectation in \eqref{eq:offline-majority-bridge} yields
\[
  \E_{\mathrm{piv}}[\dist(F^*,Y_Q)\mid Q]
  \le4\err_{H_Q}(F^*)
  =4|B_Q(F^*)|.
\]
Finally, \cref{lem:majority-half-slack} gives
$|B_Q(F^*)|\le\Delta_Q(F^*)/2$.
\end{proof}

The factor $2$ is an expected local constant because \Pivot{} is randomized.
Accordingly, we combine its expectation with the expectation over sampled
triples and then apply Markov's inequality; we do not invoke a deterministic
local-solver theorem verbatim.

\subsection{Existence of a good offline candidate}
\label{subsec:offline-candidate-existence}

Set
\begin{equation}
  \alpha:=\frac{1}{2\binom32+1}=\frac17.
  \label{eq:offline-alpha}
\end{equation}
Let
\[
  \mathcal G:=\{i\in[m]:
      \cost_{\Pp}(P_i)\le(2-\alpha)\OPT\}.
\]

\begin{lemma}[A good candidate is sampled]
\label{lem:offline-good-candidate}
With high probability, the candidate family $\Y$ contains a partition $C$
such that
\begin{equation}
  \cost_{\Pp}(C)
  \le\left(\frac{13}{7}+18\delta\right)\OPT.
  \label{eq:offline-pre-evaluation-ratio}
\end{equation}
\end{lemma}

\begin{proof}
Again there are two cases.

\emph{Many good inputs.}
If $|\mathcal G|\ge\delta m$, one input sample hits $\mathcal G$ with
probability at least $\delta$.  Thus
$q_{\mathrm{in}}=\Theta(\delta^{-1}\log n)$ independent samples contain a
member of $\mathcal G$ with high probability.  Its cost is
\[
  (2-\alpha)\OPT=\frac{13}{7}\OPT,
\]
which is stronger than \eqref{eq:offline-pre-evaluation-ratio}.

\emph{Few good inputs.}
Suppose $|\mathcal G|<\delta m$.  The universal averaging lemma with $r=3$
and \cref{lem:offline-local-C2} give
\begin{align*}
  \E_{Q,\mathrm{piv}}[\dist(F^*,Y_Q)]
  &\le2\E_Q[\Delta_Q(F^*)]\\
  &\le6(\alpha+2\delta)\OPT.
\end{align*}
By Markov's inequality,
\begin{equation}
  \Prb\!\left[
    \dist(F^*,Y_Q)
    \le6(\alpha+3\delta)\OPT
  \right]
  \ge\frac{\delta}{\alpha+3\delta}
  =\Theta(\delta).
  \label{eq:offline-local-success-probability}
\end{equation}
On this event, the metric triangle inequality implies
\begin{align*}
  \cost_{\Pp}(Y_Q)
  &\le\cost_{\Pp}(F^*)+\dist(F^*,Y_Q)\\
  &\le\bigl(1+6\alpha+18\delta\bigr)\OPT\\
  &=\left(\frac{13}{7}+18\delta\right)\OPT.
\end{align*}
By \eqref{eq:offline-local-success-probability},
$q_{\mathrm{loc}}=\Theta(\delta^{-1}\log n)$ independent trials include
such a candidate with high probability.
\end{proof}

\subsection{Selecting the best candidate from a fresh sample}
\label{subsec:offline-evaluation}

For completeness, we prove the candidate-evaluation guarantee used in the
last line of \cref{alg:offline-global}.  It is the standard metric-median
sampling argument of Indyk~\cite{Indyk1999,Indyk2001}.

\begin{lemma}[Evaluation sample]
\label{lem:evaluation-sample}
Let $\delta\in(0,1]$, and let $\Y$ be any finite candidate family independent of a uniform evaluation
sample $S$ of input partitions.  If
\[
  |S|=\Theta\!\left(
    \delta^{-2}\bigl(\log|\Y|+\log(1/\zeta')\bigr)
  \right),
\]
then, with probability at least $1-\zeta'$, the empirical-cost minimizer
$\widehat C$ satisfies
\[
  \cost_{\Pp}(\widehat C)
  \le(1+\delta)\min_{C\in\Y}\cost_{\Pp}(C).
\]
\end{lemma}

\begin{proof}
Condition on $\Y$, and let $A$ be its minimum-true-cost candidate.  Fix a
candidate $B$ with
$\cost_{\Pp}(B)>(1+\delta)\cost_{\Pp}(A)$.  For a uniform input $P$, define
\[
  X_P:=\dist(B,P)-\dist(A,P).
\]
Then
\[
  \E[X_P]=\cost_{\Pp}(B)-\cost_{\Pp}(A).
\]
The strict cost inequality implies $\dist(A,B)>0$, so the divisions below are
well defined.
The reverse triangle inequality gives $|X_P|\le\dist(A,B)$, while the triangle
inequality averaged over $P$ gives
\[
  \dist(A,B)
  \le\cost_{\Pp}(A)+\cost_{\Pp}(B).
\]
Consequently,
\[
  \frac{\E[X_P]}{\dist(A,B)}
  \ge
  \frac{\cost_{\Pp}(B)-\cost_{\Pp}(A)}
       {\cost_{\Pp}(B)+\cost_{\Pp}(A)}
  >\frac{\delta}{2+\delta}.
\]
Hoeffding's inequality, after scaling by $\dist(A,B)$, shows that $B$ has
empirical cost at most that of $A$ with probability
$\exp(-\Omega(\delta^2|S|))$.  A union bound over all $B\in\Y$ proves that no
candidate whose true cost exceeds $(1+\delta)\cost_{\Pp}(A)$ can defeat $A$
on the evaluation sample.
\end{proof}

In \cref{alg:offline-global}, $|\Y|=O(\delta^{-1}\log n)$ and
$s=\Theta(\delta^{-2}\log(n/\delta))$.  Thus
$s=\Omega(\delta^{-2}(\log n+\log|\Y|))$, so the lemma holds with
polynomially small failure probability.

\subsection{Proof of the offline theorem and resource bounds}
\label{subsec:offline-final}

Combining \cref{lem:offline-good-candidate,lem:evaluation-sample}, with high
probability the returned solution satisfies
\begin{equation}
  \cost_{\Pp}(\widehat C)
  \le(1+\delta)\left(\frac{13}{7}+18\delta\right)\OPT.
  \label{eq:offline-final-internal}
\end{equation}
For the accuracy $\eps$ in \cref{thm:offline-main}, choose
\begin{equation}
  \delta:=\min\left\{\frac1{14},\frac{\eps}{40}\right\}.
  \label{eq:offline-delta-choice}
\end{equation}
Since $\delta\le1/14$, elementary expansion of
\eqref{eq:offline-final-internal} shows
\[
  (1+\delta)\left(\frac{13}{7}+18\delta\right)
  \le\frac{13}{7}+40\delta
  \le\frac{13}{7}+\eps.
\]
This proves the approximation statement.

By \cref{lem:implicit-pivot-linear}, the local trials use
\[
  O(q_{\mathrm{loc}}n)
  =O(\delta^{-1}n\log n)
\]
time.  The candidate family has
$|\Y|=O(\delta^{-1}\log n)$ members, and the evaluation sample has
$s=O(\delta^{-2}\log(n/\delta))$ members.  By \cref{lem:linear-distance}, all
candidate evaluations take
\begin{equation}
  O\!\left(|\Y|sn\right)
  =O\!\left(\delta^{-3}n\log n\log\frac{n}{\delta}\right)
  \label{eq:offline-runtime}
\end{equation}
time.  Reading and storing the sampled input candidates costs
$O(q_{\mathrm{in}}n)$, and loading the evaluation arrays costs $O(sn)$; both
terms are dominated by \eqref{eq:offline-runtime}.  The candidate arrays occupy
$O(\delta^{-1}n\log n)$ words; the evaluation inputs can be scanned one at a
time.  Hence the working space is $\wtO_\delta(n)$.

Finally, the algorithm reads
\[
  q_{\mathrm{in}}+3q_{\mathrm{loc}}+s
  =O\!\left(\delta^{-2}\log\frac{n}{\delta}\right)
\]
input label arrays.  With the constant choice in
\eqref{eq:offline-delta-choice}, these bounds are
$\wtO_{\eps}(n)$ time and working space and $O_{\eps}(\log n)$ inspected inputs, as
claimed in \cref{thm:offline-main}.

\begin{remark}[Why the majority graph is only local]
The original consensus objective is weighted, with a different confidence
margin for every pair.  The unweighted graph $H_Q$ is formed only from a
three-input local sample and is used to construct a candidate close to the
unknown optimum.  The final candidate is evaluated against fresh samples from
the original collection.  Thus the proof does not assert that unweighted
majority error is a multiplicative surrogate for the full weighted objective.
\end{remark}

% =============================================================================
% =============================================================================
\section{Fully Dynamic Consensus Clustering}
\label{sec:dynamic}
% =============================================================================

We now prove \cref{thm:dynamic-main,thm:adaptive-main}.  The main algorithmic
ingredient is the near-linear-time snapshot algorithm of
\cref{sec:offline}.  A complete recomputation is affordable because it works
with label arrays and never constructs the dense pairwise instance.  The
remaining issues are probabilistic and representational: we must obtain a
guarantee over an entire update sequence, retain enough input information to
handle deletions, and distinguish a fixed sample against an oblivious stream
from fresh sampling against an adaptive adversary.

\subsection{Update model and the meaning of the resource bounds}
\label{subsec:dynamic-resources-model}

We use the vertex-update model of \cref{subsec:dynamic-model}, with fixed
$m$, at most $N$ active vertices, and at most $T$ updates.  The first
construction uses sampled coordinates and an oblivious stream; the adaptive
extension uses the full-current-input model.  Both include an optional initial
snapshot at time $0$.  Deleting a vertex restricts the input partitions and
never changes the relation between two survivors.

The space bound for the oblivious construction includes all sampled labels
retained for active vertices.  Thus a future deletion needs no new input
queries.  If a complete incoming vector must be scanned, the additional
$O(m)$ ingestion cost is charged only to insertions.  For the adaptive
extension, the full current input database is a separate, explicitly charged
resource.

Labels and vertex identifiers need not be consecutive integers.  They are
assumed to occupy one word.  The word size also accommodates the sample
indices, pointers, and integer score sums used below; in particular, a sum
of $s$ distances is at most $s\binom{N}{2}$.  We rank-compact the labels of
the current snapshot before using the linear-time primitives of
\cref{sec:offline}.  The label arrays output by our algorithm use ranks in
$[n_t]$.  These output labels are snapshot identifiers, and we impose no
recourse guarantee on their changes.

\subsection{A snapshot theorem with an explicit failure parameter}
\label{subsec:snapshot-confidence}

We first make the confidence dependence of the offline result explicit.
This will allow a direct union bound over an update sequence whose length
can be much larger than its maximum number of active vertices.  All
logarithms in this section are natural unless a base is displayed.

\begin{lemma}[Snapshot approximation at confidence $1-\beta$]
\label{lem:snapshot-confidence}
Let $\delta\in(0,1/14]$ and $\beta\in(0,1)$.  For a fixed instance with
$m\ge3$ and $n\ge2$, set
\begin{equation}
  q_\beta:=\left\lceil\delta^{-1}\log\frac2\beta\right\rceil,
  \qquad
  s_\beta:=\left\lceil18\delta^{-2}
                  \log\frac{4q_\beta}{\beta}\right\rceil.
  \label{eq:snapshot-explicit-sizes}
\end{equation}
Draw $q_\beta$ independent uniform input indices and $q_\beta$ independent
uniform three-element subsets of $[m]$.  Form one candidate from each sampled
input and one independent \Pivot{} candidate from each sampled triple.
Select the minimum empirical-cost candidate using an independent multiset
of $s_\beta$ uniform input indices.  With probability at least $1-\beta$,
the output $\widehat C$ satisfies
\begin{equation}
  \cost_{\Pp}(\widehat C)
  \le(1+\delta)\left(\frac{13}{7}+18\delta\right)\OPT.
  \label{eq:snapshot-confidence-ratio}
\end{equation}
More generally, the same assertion holds for any $q\ge q_\beta$ input
candidates and $q$ local candidates provided the evaluation sample has size
at least $\lceil18\delta^{-2}\log(4q/\beta)\rceil$.
\end{lemma}

\begin{proof}
Set $\alpha=1/7$ and call an input index good if its partition has cost at
most $(2-\alpha)\OPT$.  If at least a $\delta$ fraction of the indices
are good, the probability that all $q$ input samples miss them is at most
\[
  (1-\delta)^q\le e^{-\delta q}\le\frac\beta2.
\]
Otherwise, the calculation in
\cref{lem:offline-local-C2,lem:offline-good-candidate} shows that one local
trial produces a candidate of cost at most
$(13/7+18\delta)\OPT$ with probability at least
\[
  \frac{\delta}{\alpha+3\delta}\ge\delta.
\]
The last inequality follows from
$\alpha+3\delta\le1/7+3/14=5/14<1$.  Hence all $q$ independent local
trials fail with probability at most $e^{-\delta q}\le\beta/2$.
Only one of these two cases is needed for a fixed instance.  It follows
that the probability of failing to generate a candidate of the stated
cost is at most $\beta/2$, rather than the sum of two such terms.

For completeness, we give the numerical confidence bound for candidate
evaluation.  Condition on the candidate family $\Y$, of size at most
$2q$, and let $A\in\Y$ minimize true cost.  For a candidate $B$ whose
cost exceeds $(1+\delta)\cost_{\Pp}(A)$, put
$D:=\dist(A,B)>0$.  The normalized cost-difference random variable
\[
  Z_P:=\frac{\dist(B,P)-\dist(A,P)}D
\]
belongs to $[-1,1]$, and the argument in
\cref{lem:evaluation-sample} gives
\[
  \E[Z_P]>\frac{\delta}{2+\delta}\ge\frac\delta3.
\]
Hoeffding's inequality therefore bounds the probability that $B$ has
empirical cost at most that of $A$ by
\[
  \exp\!\left(-\frac{s\delta^2}{18}\right).
\]
A union bound over at most $2q$ candidates makes the probability of an
incorrect empirical choice at most
\[
  2q\exp\!\left(-\frac{s\delta^2}{18}\right)
  \le\frac\beta2.
\]
Together with the candidate-generation bound, this proves the lemma.

There is no exception when $\OPT=0$.  In that case all input partitions
coincide with the zero-cost optimum, so every input index is good and the
first case applies.  The evaluation proof considers only candidates with
strictly larger cost, for which $D>0$ remains valid.
\end{proof}

The lemma concerns a fixed snapshot and does not assume that its optimum is
related to the optimum of any earlier snapshot.  This independence from
temporal monotonicity is the reason arbitrary deletions cause no difficulty
for the approximation analysis.

\subsection{Fixed samples and a complete rebuild after each update}
\label{subsec:dynamic-algorithm}

Fix the accuracy $\eps>0$ and total failure probability $\zeta\in(0,1)$,
and define
\begin{equation}
  \delta:=\min\left\{\frac1{14},\frac\eps{40}\right\},
  \qquad
  L:=\log\frac{8(N+1)(T+1)}{\delta\zeta},
  \label{eq:dynamic-rebuild-parameters}
\end{equation}
and
\begin{equation}
  q:=\left\lceil2\delta^{-1}L\right\rceil,
  \qquad
  s:=\left\lceil72\delta^{-2}L\right\rceil.
  \label{eq:dynamic-sample-sizes}
\end{equation}
We first describe the algorithm for $m\ge3$.

At initialization, draw independent uniform input indices
$I_1,\ldots,I_q\in[m]$, independent uniform triples
$Q_1,\ldots,Q_q\in\binom{[m]}3$, and an independent evaluation multiset
$S=(J_1,\ldots,J_s)$ of uniform input indices.  Retain these samples for
the entire execution.  Let
\[
  R:=\{I_a:a\in[q]\}\cup
      \bigcup_{a\in[q]}Q_a\cup
      \{J_b:b\in[s]\}.
\]
The set $R$ contains at most $4q+s$ distinct indices.  We store the
labels $c^i(v)$ only for active vertices $v$ and indices $i\in R$.
Repeated indices may share a stored array, but their multiplicities in
sampling and empirical scoring are preserved.

After initialization and after every update, form the current candidate family
\begin{equation}
  \Y_t:=\{P_{I_a}^t:a\in[q]\}
      \cup\{Y_{a,t}:a\in[q]\},
  \qquad
  Y_{a,t}:=\Pivot(H_{Q_a}^t).
  \label{eq:dynamic-candidates}
\end{equation}
Every \Pivot{} execution uses fresh randomness, independent of the samples,
of the other executions, and of all previous executions.  Crucially,
candidate generation does not use the evaluation sample $S$.  Return a
partition minimizing
\begin{equation}
  \widehat{\cost}_t(C)
    :=\frac1s\sum_{b=1}^s\dist_{V_t}(C,P_{J_b}^t),
  \qquad C\in\Y_t.
  \label{eq:dynamic-empirical-cost}
\end{equation}
Ties can be broken by a fixed rule.  One may compare the integer score sums
without performing any divisions.

\begin{paperalgorithm}[htbp]
\centering
\fbox{\begin{minipage}{0.93\textwidth}
\textbf{Fully dynamic consensus with persistent input samples.}

\begin{enumerate}[(1)]
  \item \textbf{Initialize.} Sample $(I_a)_{a\in[q]}$,
        $(Q_a)_{a\in[q]}$, and $S$ independently, with the sizes in
        \eqref{eq:dynamic-sample-sizes}.  Store their distinct index set $R$
        and the labels of every initially active vertex in these inputs.
  \item \textbf{Insert $v$.} Read $c^i(v)$ for every $i\in R$ and add a
        row for $v$ to the retained input representation.
  \item \textbf{Delete $v$.} Remove the row of $v$ from all retained input
        arrays.  The retained labels of surviving vertices stay unchanged.
  \item \textbf{Prepare the snapshot.} Rank-compact each retained input's
        labels on the current active set.  If $n_t\le1$, output the unique
        partition and finish this update.
  \item \textbf{Rebuild candidates.} Add the $q$ sampled input restrictions
        to $\Y_t$.  For every sampled triple, execute the implicit
        \Pivot{} procedure of \cref{alg:implicit-pivot} with fresh coins
        and add its output to $\Y_t$.
  \item \textbf{Evaluate and output.} Compute the distances between all
        candidates and all sampled evaluation inputs using
        \cref{lem:linear-distance}.  Output a label array for a minimizer
        of \eqref{eq:dynamic-empirical-cost}.
\end{enumerate}
\end{minipage}}
\caption{The oblivious-stream algorithm.  Steps (2) and (3) are alternative
update operations; steps (4)--(6) follow either operation.  Input samples
persist, while the local clusterings are rebuilt.  Steps (4)--(6) also run
once after initialization to produce the optional initial output.}
\label{alg:dynamic-rebuild}
\end{paperalgorithm}

At a snapshot with no vertices, the unique partition is the empty family;
with one vertex, it is the singleton partition.  We still retain the
sampled labels of an inserted vertex even when the current output is
trivial, because those labels will be needed at later snapshots.  If
$m\le2$, sampling is unnecessary: following $P_1^t$ is exact at every
snapshot by the metric triangle inequality.  Storing that one input and
materializing its restriction fits all the bounds below.

\subsection{A simultaneous guarantee under arbitrary oblivious deletions}
\label{subsec:dynamic-correctness}

The next lemma formalizes the use of persistent samples.  Reusing a sample
across time does not make the candidate family and the evaluation sample
independent after conditioning on previous outputs.  No such conditioning
is used in the oblivious-stream proof.

\begin{lemma}[Marginal snapshot law]
\label{lem:dynamic-marginal-law}
Fix an oblivious stream and a snapshot $t$.  The random candidates and
evaluation sample used by \cref{alg:dynamic-rebuild} at time $t$ have the
same joint distribution as a fresh execution of the sampling algorithm
in \cref{lem:snapshot-confidence} on the fixed instance $\Pp^t$, with
$q$ input samples, $q$ local samples, and $s$ evaluation samples.
\end{lemma}

\begin{proof}
Because the stream is oblivious, $V_t$ and each partition $P_i^t$ are fixed
objects independent of all random choices.  The indices $I_a$ remain
independent uniform elements of $[m]$, and the $Q_a$ remain independent
uniform elements of $\binom{[m]}3$.  The stored labels reconstruct exactly
the restrictions of those inputs to $V_t$, regardless of the insertions
and deletions that led to this active set.  Fresh, independent pivot coins
therefore give precisely the candidate-generation distribution prescribed
by the snapshot algorithm.

The evaluation indices were drawn independently of the candidate indices,
the triple indices, and every pivot execution.  Thus they are independent
of $\Y_t$, which is a function only of the fixed instance and those latter
random variables.  An accidental equality between an evaluation index and
a candidate index does not invalidate independence of their sampling
experiments.  This establishes the full joint distribution claimed.
\end{proof}

\begin{proof}[Approximation part of \cref{thm:dynamic-main}]
Let $\beta:=\zeta/(T+1)$.  The choices in
\eqref{eq:dynamic-parameters} and \eqref{eq:dynamic-sample-sizes} satisfy the sample
requirements of \cref{lem:snapshot-confidence}.  Indeed,
$L\ge\log(2/\beta)$, and therefore
$q\ge\delta^{-1}\log(2/\beta)$.  Also
$q\le3\delta^{-1}L$ since $\delta^{-1}L\ge1$.  It follows that
\begin{align*}
  \log\frac{4q}{\beta}
  &\le\log\frac{12\delta^{-1}L(T+1)}{\zeta}\\
  &=L+\log\frac{3L}{2(N+1)}
   \le 2L.
\end{align*}
The last inequality uses $N+1\ge1$, $L\ge\log112>1$, and
$\log(3L/2)\le L$ in this range.  Thus
$s\ge18\delta^{-2}\log(4q/\beta)$, with room to spare.

By \cref{lem:dynamic-marginal-law,lem:snapshot-confidence}, for each fixed
$t$ the probability that the output violates
\[
  \cost_t(C_t)
  \le(1+\delta)\left(\frac{13}{7}+18\delta\right)\OPT_t
\]
is at most $\beta$.  There are at most $T+1$ snapshots, so a union bound
gives this inequality simultaneously for all of them with probability at
least $1-(T+1)\beta=1-\zeta$.  The failure events at different times may
be dependent; independence is unnecessary for this union bound.

Finally, since $\delta\le1/14$,
\[
  (1+\delta)\left(\frac{13}{7}+18\delta\right)
  =\frac{13}{7}+\frac{139}{7}\delta+18\delta^2
  \le\frac{13}{7}+40\delta
  \le\frac{13}{7}+\eps.
\]
Trivial snapshots and the cases $m\le2$ are exact, completing the
approximation proof.
\end{proof}

The relevant horizon in this proof is $T$, rather than $N$.  For example,
a stream may repeatedly insert and delete vertices while keeping only a
small active set.  Substituting an inverse-polynomial failure probability
in $N$ without accounting for the number of updates would not by itself
establish a guarantee over such a stream.

\subsection{Worst-case update time and total maintained space}
\label{subsec:dynamic-implementation}

We next give the implementation details needed for the resource claims in
\cref{thm:dynamic-main}.  A uniform triple can be sampled in constant time
by choosing successive ranks among the remaining indices and skipping the
at most two indices already chosen; this avoids an unbounded rejection loop.
At initialization, sorting the $4q+s$ sampled indices constructs $R$ and
the references from each sampling position to its stored array.  This takes
$O((q+s)\log(q+s))=O(\delta^{-3}L^2)$ time, within the stated bound even
when the initial set is empty.  The subsequent initial rebuild has the same
bound as an update, with $n_0$ in place of $n_t$.

Put $k:=|R|\le4q+s$.  Maintain a dense ordering
of the active vertices and one raw-label array for each sampled input.
A dictionary maps an external vertex identifier to its current row.
An insertion reads and stores $k$ labels.  A deletion finds its row and
removes it; moving the last row into the vacated position preserves a
dense representation.  A deterministic balanced dictionary suffices for
the $O(\log(n_t+2))$ identifier operations.

For a fully worst-case bound, the label arrays can be freshly allocated
and copied during each rebuild.  The previous snapshot has at most
$n_t+1$ vertices, so even during a deletion the old and new arrays together
occupy $O(k(n_t+1))$ words.  This avoids relying on an amortized resizing
bound or retaining an $N$-sized allocation when few vertices remain.
The dictionary may similarly be reconstructed during the snapshot scan;
its $O((n_t+1)\log(n_t+2))$ construction cost is dominated below.

The retained raw labels must preserve their external meaning, because a
later insertion may use one of those labels.  We therefore rank-compact
into temporary arrays instead of replacing the stored raw labels.
For every retained input, comparison-sort the pairs consisting of raw
label and vertex row, and assign equal labels equal ranks in $[n_t]$.
The total compaction time is
\begin{equation}
  O\!\left(k(n_t+1)\log(n_t+2)\right),
  \label{eq:dynamic-compaction-cost}
\end{equation}
using $O(k(n_t+1))$ words.  Once ranks are available, the bucket
constructions in \cref{lem:implicit-pivot-linear} and the pair-key sorts
in \cref{lem:linear-distance} use counting or radix sort with universe
size $n_t$ and therefore take linear time per execution.  Cluster labels
of newly produced candidates are already ranks in $[n_t]$.

There are $q$ local trials, so generating their output arrays takes
$O(q(n_t+1))$ time and space.  There are at most $2q$ candidates and $s$
evaluation positions.  All scores can consequently be computed exactly
in
\begin{equation}
  O\!\left(qs(n_t+1)\right)
  =O\!\left(\delta^{-3}(n_t+1)L^2\right)
  \label{eq:dynamic-evaluation-time}
\end{equation}
time.  This uses the contingency-table distance identity, rather than an
enumeration of the $\binom{n_t}{2}$ vertex pairs.  It is unnecessary to
store a table of distances for all candidate--evaluation pairs: maintain
one running integer score per candidate and reuse the linear-size
distance-computation workspace.

Because $k=O(\delta^{-2}L)$ and $L\ge\log(N+1)$, the compaction cost
in \eqref{eq:dynamic-compaction-cost}, all array copies, and candidate
generation are dominated by \eqref{eq:dynamic-evaluation-time}.  Writing
the selected partition as a full length-$n_t$ label array is also included.
The bound is worst-case for each update, under the constant-time uniform
sampling convention of our word-RAM model.  It holds for every outcome of
the random samples and pivot choices: the linear-time pivot implementation
does not rely on an expected running-time bound.

The raw and compacted sampled labels require $O(k(n_t+1))$ words; the
candidate arrays require $O(q(n_t+1))$; and the pivot and scoring
workspaces require $O(n_t+q+1)$.  The sample indices themselves use
$O(q+s)$ words.  The total is therefore
\begin{equation}
  O\!\left((q+s)(n_t+1)\right)
  =O\!\left(\delta^{-2}(n_t+1)L\right)
  \label{eq:dynamic-total-space}
\end{equation}
words.  This accounts for all input labels retained by the oblivious
algorithm.  A sampled input index requires a word capable of storing an
element of $[m]$; the bound does not assert that identifiers have a bit
length independent of $m$.

Together with the preceding approximation proof, these estimates complete
\cref{thm:dynamic-main}.  When $\eps$ is constant, $T$ is polynomially
bounded in $N+1$, and $\zeta$ is inverse polynomial in $N+1$, they simplify
to $O((n_t+1)\log^2(N+2))$ worst-case update time and
$O((n_t+1)\log(N+2))$ maintained words.  Without these parameter
restrictions, \eqref{eq:dynamic-parameters} retains the dependence on the
actual update horizon and confidence.

\subsection{Adaptive updates with fresh access to the full input}
\label{subsec:dynamic-adaptive}

Persistent samples provide a compact input representation, but their
independence from the current instance was essential in
\cref{lem:dynamic-marginal-law}.  If future updates may depend on past
outputs, the active instance can become correlated with the retained
samples.  We do not claim that the fixed-sample algorithm has the same
guarantee in this adaptive model.

There is a direct adaptive variant under a stronger access assumption.
Assume that, after every update, the algorithm can obtain the label array
of any chosen current input $P_i^t$ in $O(n_t+1)$ time.  The full current
family may be maintained by an external database, or it may be stored
explicitly by the algorithm.  After the adversary commits the current
update, draw entirely fresh input indices, fresh triple indices, fresh
evaluation indices, and fresh pivot coins.  Apply the snapshot algorithm
with the same $q,s$ as in \eqref{eq:dynamic-sample-sizes}.  Discard the
samples when their output has been produced.  The candidate and evaluation
samples are mutually independent, as before.

\begin{proof}[Proof of \cref{thm:adaptive-main}]
Allow the adversary to choose update $t$ as an arbitrary function of the
past transcript and its own random bits, but require that it commit the
update before the algorithm draws its fresh random bits for snapshot $t$.
Condition on this transcript and the completed current update.  Under the
conditioning, the current instance $\Pp^t$ is fixed, while all random
choices used by the snapshot solver are fresh and independent of that
instance.  Hence \cref{lem:snapshot-confidence} applies conditionally and
gives
\[
  \Prb\!\left[
    \cost_t(C_t)>
    \left(\frac{13}{7}+\eps\right)\OPT_t
    \;\middle|\;\text{past transcript and current update}
  \right]
  \le\frac{\zeta}{T+1}.
\]
Taking expectations preserves this bound unconditionally.  Summing over
at most $T+1$ snapshots shows that the probability of any failure is at
most $\zeta$.  This argument permits the adversary to observe all earlier
output partitions; it does not require the sequence to have been fixed
at initialization.

The implementation within a snapshot is identical to the oblivious
algorithm: load at most $4q+s$ current input arrays, compact their labels,
generate the local candidates, and score them.  Thus the clustering
computation takes
$O(\delta^{-3}(n_t+1)L^2)$ worst-case time and
$O(\delta^{-2}(n_t+1)L)$ working words after the current full input is
available.  These working words include the sampled arrays loaded for
this snapshot.

If the full input is maintained locally, its additional storage is
$O(m(n_t+1))$ words.  One simple representation keeps a length-$m$ raw
label vector in each active vertex record and links the records in an
active list.  Reading any chosen input coordinate over the active list
takes $O(n_t+1)$ time.  Inserting or deleting a record takes $O(m)$ time
in addition to identifier-dictionary operations, which are dominated by
the clustering computation above.  Thus a fully local implementation
has total update time
\[
  O\!\left(m+\delta^{-3}(n_t+1)L^2\right)
\]
and total space
\[
  O\!\left(m(n_t+1)+\delta^{-2}(n_t+1)L\right).
\]
If an external input store supplies current arrays, these full-input
storage and maintenance terms belong to that store, and the stated
clustering time and workspace describe the algorithm's own computation.
This completes the proof.
\end{proof}

The extra access assumption is substantive.  Fresh sampling can request
an input index that has never previously been inspected, so its labels
on all surviving vertices must be recoverable.  The compact retained
sample of \cref{thm:dynamic-main} does not provide that capability by
itself.

\subsection{Relationship to the online result and to dynamic agreement}
\label{subsec:dynamic-scope}

Our online and dynamic theorems optimize different resources.  The online
algorithm follows persistent local solutions and uses history scores to
bound the total number of identifier changes of every vertex.  The
dynamic algorithm recomputes local solutions and selects a fresh snapshot
output, thereby allowing the current optimum to decrease after a deletion.
It may change the output labels of $\Theta(n_t)$ surviving vertices in one
update.  No polylogarithmic per-vertex recourse statement follows from its
update-time or approximation bound.

The dynamic theorem is a consequence of the near-linear snapshot solver
together with persistent sampling and explicit access accounting.  It
does not maintain a dense weighted correlation-clustering instance under
individual edge updates.  In particular, it uses no fully dynamic
correlation-clustering algorithm as a black box.

Dynamic Agreement of Cohen-Addad et al.~\cite{CLMP2024Dynamic} addresses
unweighted correlation clustering under a different update and access
model.  Its sublinear-update analysis allows adversarial vertex insertions
and assumes random vertex deletions.  Our oblivious theorem permits any
predetermined sequence of deletions, but pays near-linear time to rebuild
the current output.  These statements concern different objectives and
different resource guarantees.  Substituting a dynamic routine for
\Agree{} inside the online theorem would require the approximation and
persistent-identifier guarantees used in
\cref{thm:agree-blackbox}; a fast dynamic update bound alone does not supply
that interface.

Finally, when an explicit output partition is requested after every
update, writing its $n_t$ labels already costs $\Omega(n_t)$ operations.
The dynamic computation is within polylogarithmic factors of this output
cost under the sampled-coordinate assumptions.  This observation is only an
output-size comparison: it is not a lower bound for an implicit-output or
query-only dynamic data structure.

% =============================================================================
\section{Conclusion and Open Problems}
\label{sec:conclusion}
% =============================================================================

The common explanation behind our results is that the factor-$2$ benchmark
for metric medians can be crossed when a constant-size
sample supports a useful candidate not already present among the inputs.  For
partition distance, three inputs suffice: their majority relation has error at
most half of the total triangle slack.  Offline, an implicit linear-time
\Pivot{} execution turns this relation into a $13/7+\eps$ approximation while
reading only polylogarithmically many inputs.  Online, a persistent deterministic
majority solver and monotone history scores make the same phenomenon stable
throughout the vertex-arrival sequence.

For fully dynamic inputs, we use the same snapshot solver after every vertex
insertion or deletion.  Fixing the sampled input indices gives a compact
representation against oblivious updates, while fresh sampling gives an
adaptive guarantee when all current inputs remain accessible.  A union bound
over the update horizon yields a simultaneous $13/7+\eps$ approximation.
The implementation includes arbitrary-label compaction, exact evaluation on
the sampled inputs, and materialization of the complete output in
$O_{\eps}((n_t+1)L^2)$ worst-case time.  The dynamic theorem follows from
fast recomputation and therefore controls computation without controlling
the lifetime recourse of individual vertices.

Several directions remain open.  First, the Online Agreement constant $\eta$
is not explicit in its published analysis.  A sharper, explicit persistent
solver for three-partition majority graphs would immediately improve the online
competitive factor.  Second, our work leaves open a deterministic sub-$2$
guarantee at every prefix of an oblivious vertex-arrival sequence with
polylogarithmic persistent-identifier recourse; our randomized proof uses
initial sampling as a hitting family over all prefixes.  Third, the online
algorithm maintains exact candidate costs and does not optimize update time.
Combining history-stable selection with sublinear evaluation is a natural next
step.  In the fully dynamic model, achieving a sub-$2$ approximation together
with both near-linear update time and small per-vertex recourse is a separate
challenge, especially when deletions sharply reduce the optimum.  Another
question is whether an adaptive algorithm can retain only a small sampled
input representation while obtaining the guarantees currently proved using
fresh access to all inputs.  Finally, the offline and dynamic constants could
improve if a better fast local solver is implemented on the implicit
triple-majority representation, while preserving the worst-case near-linear
access bound.

\newpage
\ifusebiblatex
  \printbibliography
\else
  \bibliographystyle{alpha}
  \bibliography{ref}
\fi

\end{document}
